\documentclass[11pt]{article}

\usepackage[a4paper,margin=2.6cm]{geometry}
\usepackage{amsmath}
\usepackage{amssymb}
\usepackage{amsthm}
\usepackage{booktabs}
\usepackage{xcolor}
\usepackage{enumitem}
\usepackage{longtable}
\usepackage[colorlinks=true,linkcolor=blue!55!black,urlcolor=blue!55!black]{hyperref}

\newsavebox{\kbox}
\newenvironment{keypoint}
  {\begin{lrbox}{\kbox}\begin{minipage}{0.93\textwidth}\vspace{3pt}}
  {\vspace{3pt}\end{minipage}\end{lrbox}\begin{center}\fbox{\usebox{\kbox}}\end{center}}

% Every central equation gets the same three-part gloss.
\newcommand{\says}[1]{\smallskip\noindent\textbf{What it says.} #1\smallskip}
\newcommand{\why}[1]{\smallskip\noindent\textbf{Why we model it this way.} #1\smallskip}
\newcommand{\from}[1]{\smallskip\noindent\textbf{Where it comes from.} #1\smallskip}
\newcommand{\checked}[1]{\smallskip\noindent\textbf{How this was checked.} #1\smallskip}

\theoremstyle{plain}
\newtheorem{theorem}{Theorem}
\newtheorem{proposition}{Proposition}
\newtheorem{lemma}{Lemma}
\newtheorem{corollary}{Corollary}
\theoremstyle{remark}
\newtheorem{remark}{Remark}
\theoremstyle{definition}
\newtheorem{assumption}{Assumption}
\newtheorem{definition}{Definition}

\setlist[itemize]{itemsep=2pt,topsep=4pt}
\setlist[enumerate]{itemsep=3pt,topsep=4pt}

\title{\bfseries Multinational Ownership, Market Power\\ and the Measurement of Foreign Profits\\[8pt]
\large A general equilibrium model with large firms and entry:\\
what each equation says, where it comes from,\\
and proofs that the model has exactly one solution}
\author{Sebasti\'an Vel\'asquez Palacios \and Christian Volpe Martincus\\[2pt]
\normalsize Inter-American Development Bank}
\author{Sebasti\'an Vel\'asquez Palacios \and Christian Volpe Martincus\\[2pt]
\normalsize Inter-American Development Bank}
\date{\today\ --- working draft, please do not circulate}

\begin{document}
\maketitle

\begin{abstract}
\noindent How much of the profit that multinational firms earn in a market belongs
to the country that sets policy toward it? We build a general-equilibrium model of
multinational production in which markups rise with an ownership group's market
share, ownership is recorded firm by firm, and the set of operating factories is an
equilibrium outcome. The model has exactly one solution under stated, checked
conditions: the market game and the entry game are proved unique, and the wage layer
is certified on a continuum of wage vectors by interval arithmetic. The central
result is an identity: the standard country-level calculation of foreign profit errs
by exactly the covariance between ownership and firm size --- measured at $+10\%$
for the United States in matched Orbis/D\&B--customs data for ten Latin American
countries, and sign-varying across owners. The model confronts six stylized facts of
multinational exporting and reproduces them, three out of sample.
\end{abstract}

\begin{keypoint}
\textbf{How this document is organised.} Every central equation is followed by three
short paragraphs, always in the same order: \emph{what it says} in ordinary words,
\emph{why we model it this way} rather than some other way, and \emph{where it comes
from} in the literature. Every symbol is defined before it is used, and all of them
are collected in Appendix~\ref{sec:notation}. Nothing is left as ``it can be
shown'': where a claim is proved, the proof is written out; where a claim is only
checked numerically, it says so in those words, and every such claim is collected in
Section~\ref{sec:notproved}.
\end{keypoint}

\tableofcontents

\newpage
\section{Introduction}
\label{sec:intro}


A group of developing countries sells goods abroad. Much of what they sell is shipped
not by local firms but by subsidiaries of large international groups whose head
offices are somewhere else --- frequently in the very country that imposes tariffs on
those goods.

We want to answer one question: \emph{how much of the profit earned in those markets
ends up belonging to the country that sets the tariff?}

The obvious way to answer it is to take total profit and multiply by the share of the
firms that the country owns. Section~\ref{sec:result} proves that this is wrong, shows the direction of
the error, and gives a formula for its exact size. Getting that formula requires a
model with three specific features, and Section~\ref{sec:three} explains why each one
is unavoidable.

\subsection{Three features the model must have}
\label{sec:three}

\begin{enumerate}
\item \textbf{Large firms whose markups rise with size.} If every firm charged the
same markup, profit would be a fixed fraction of sales, and ownership would be a pure
transfer with no consequences for anything real. Proposition~\ref{prop:irrelevance}
proves this formally: with constant markups the question above has no content. The
standard workhorse model of international trade therefore cannot be used, and that is
a matter of logic rather than preference.
\item \textbf{Firms headquartered in one country, producing in a second, selling to a
third.} This is what the data record, and it is what makes ``who owns the profit''
different from ``who produced the goods''.
\item \textbf{Firms that enter and exit.} Which factories exist must respond to wages
and to policy. Otherwise the most interesting patterns in the data are inserted by
hand rather than explained.
\end{enumerate}

\noindent Each feature is standard on its own; what is new here is the combination
together with \emph{firm-level ownership data}. The closest relatives each lack one
piece: Yang (2023) has multi-plant oligopoly and location choice but no ownership
accounting; Alviarez, Head and Mayer (2025) have oligopoly with multinational brand
ownership but no entry margin or tariff question; the quantitative multinational-production
benchmarks (Ramondo--Rodr\'iguez-Clare 2013; Tintelnot 2017) have constant or no markups,
so ownership is irrelevant in them by construction. Combining the three features creates
a technical problem --- models with multi-factory firms and entry usually admit many
possible outcomes --- and Section~\ref{part:proofs} is devoted to solving it.

\subsection{Related literature: where every piece of the model comes from}
\label{sec:provenance}

\says{Almost none of the machinery is new, and it is worth saying in two paragraphs
which parts are borrowed and from whom before the table sets it out piece by piece.}

The market structure comes from Atkeson and Burstein (2008): two layers of competition,
where goods compete closely inside a sector and loosely across sectors, which is what
makes a firm's markup depend on how big it is. The one change is that the markup is
computed at the level of the \emph{owner} rather than the factory. A group that owns
several plants in the same market internalises the competition among them, so what sets
its markup is its combined share, not each plant's --- the mechanism Yang (2023) calls
cannibalisation, and the reason ownership shows up in prices at all. With constant
markups the entire result vanishes, which is why monopolistic competition is not an
option here. Which factories exist is then an outcome rather than a list, and the
granular machinery for that is Gaubert and Itskhoki (2021): capability drawn from a
Pareto distribution, and a pool of potential parents that grows with country size, so
the concentration of parent firms in a handful of rich countries comes from a bigger
pool having a better maximum rather than from assuming that rich countries breed better
firms. Entry is settled by a theorem rather than by an assumed order of moves, which
matters because an order of moves is not a property of the model.

The multinational side comes from Ramondo and Rodr\'iguez-Clare (2013): producing away
from head office is less efficient, a friction measured for the car industry by Head and
Mayer (2019). Tintelnot (2017) supplies the geography --- shipping is counted from where
a good is \emph{made} rather than from head office, so one factory can serve many
countries, and serving each of them costs a fixed amount, which is what turns location
into a decision. Caliendo and Parro (2015) supply the input--output loop, so that firms
buy from each other at the prices prevailing where they sit. Head-office intensity is
Antr\`as (2003) and Antr\`as and Helpman (2004), but used on the \emph{capability} side
rather than the cost side: a group's head office is what makes its factories good, not a
service each factory buys at its parent's wage. Complexity therefore raises the penalty
on a stand-alone local firm and raises the return to a group's capability, but it never
puts a foreign wage into a domestic cost --- a choice that turns out to decide whether
the model has one solution, and which improves the fit besides.\footnote{Moving
head-office services from the cost side to the capability side raises the Fact~2
gradient from $+0.33$ to $+0.42$ against a target of $+0.43$, and makes the fitted
multinational productivity edge unnecessary, so four internally calibrated parameters
become three. Section~\ref{sec:hqchoice} sets out the choice and
Section~\ref{sec:contraction} the consequence for uniqueness.} The version where
head-office services are bought at the parent's wage is supported and reported
alongside. What is not borrowed from anyone: ownership is at the firm level, $\theta$,
so who owns which factory is data rather than a country-level average --- and that is
the whole point.

\noindent The table below names the source of each component and what it does here. The
two rows in bold are the contribution. Full references, with years and outlets checked
against the published articles, are in Section~\ref{sec:refs}.

\begin{center}
\small
\begin{longtable}{p{4.1cm}p{4.5cm}p{5.8cm}}
\toprule
\textbf{Component} & \textbf{Source} & \textbf{What it does here}\\
\midrule
\endfirsthead
\toprule
\textbf{Component} & \textbf{Source} & \textbf{What it does here}\\
\midrule
\endhead
Two layers of competition: goods compete closely inside a sector, loosely across
sectors &
Atkeson and Burstein (2008) &
Makes a firm's markup depend on its size. With one layer only, every firm charges the
same markup and Section~\ref{sec:result} has nothing to say.\\
\addlinespace
The markup depends on the \emph{group's} total share, not the individual factory's &
The same two-layer logic, extended here to multi-factory international groups &
The central mechanism, which we call \emph{cannibalisation}. Everything in Section~\ref{sec:result}
follows from it.\\
\addlinespace
Producing away from head office is less efficient ($\gamma_{h\ell}$) &
Ramondo and Rodr\'iguez-Clare (2013); measured for the car industry by Head and Mayer
(2019) &
Determines where a group can profitably build, and hence who becomes an international
group at all.\\
\addlinespace
Shipping counted from the \emph{factory} rather than from head office ($d_{\ell n}$)
&
Tintelnot (2017) &
Lets one factory serve many countries, and produces Fact~6 directly.\\
\addlinespace
A fixed cost of selling into each market ($F_{a,n}$) &
Tintelnot (2017) &
Turns ``which markets do I serve'' into a decision rather than an assumption.\\
\addlinespace
Head-office work matters more in complex products ($\alpha_k$) &
Antr\`as (2003); Antr\`as and Helpman (2004) &
The reason complex products and foreign ownership go together --- with a twist
explained in Section~\ref{sec:fact2}.\\
\addlinespace
Firms buy inputs from other firms, sector by sector ($\nu_k$, $\omega_{k'k}$) &
Caliendo and Parro (2015) &
Lets international firms lower local input prices --- a channel Fact~5 was hoped to
use, measured insufficient (Section~\ref{sec:fact5}).\\
\addlinespace
A few very large firms; the \emph{number} of possible firms grows with country size;
entry ranked by cost &
Gaubert and Itskhoki (2021); Gaubert, Itskhoki and Vogler (2021) &
Produces Facts~1 and 3 without assuming that rich countries breed better firms.\\
\addlinespace
Groups with several factories choosing \emph{which set} to operate &
Yang (2023) &
The market structure --- a few large groups plus many small local firms --- and the
multiplicity problem Section~\ref{part:proofs} solves.\\
\addlinespace
International firms may raise local firms' \emph{productivity} &
Javorcik (2004) &
Switched off in the baseline; tested against Fact~5 and shown insufficient at any
size (Section~\ref{sec:fact5}).\\
\addlinespace
International presence may lower local firms' \emph{cost of reaching export markets} &
Aitken, Hanson and Harrison (1997) --- the classic evidence that exporting spills over
from multinationals to neighbours &
Switched off in the baseline; tested against Fact~5: it fully repairs the entry margin
and the value response is still $-80\%$ (Section~\ref{sec:fact5}).\\
\addlinespace
Outward-rounded interval arithmetic; contraction of the Hilbert metric &
Moore (1966); Birkhoff (1957) &
The continuum certificate for wage uniqueness: every wage vector in a box at once,
no between-points gap (Section~\ref{sec:notproved}).\\
\addlinespace
Cross-border ownership and the optimal tariff &
Brecher and Bhagwati (1981); Blanchard (2007, 2010); Blanchard and Matschke (2015);
Blanchard, Bown and Johnson (2026); Itskhoki and Mukhin (2025) &
The existing literature. Every one of these uses a \emph{country-level} ownership
share --- exactly what this model shows is not enough.\\
\addlinespace
Comparative-statics tools used in the proofs &
Milgrom and Shannon (1994); Novshek (1985) &
The standard results behind the two steps of Theorem~\ref{thm:entry}. The proof given
here is self-contained.\\
\midrule
\textbf{Ownership recorded firm by firm ($\theta_{gn}$), and the income formula it
implies} &
\textbf{This model} &
\textbf{Proposition~\ref{prop:main}: the correct adjustment involves market
concentration, not an average ownership share.}\\
\addlinespace
\textbf{A proof that entry has a unique outcome even when groups compete with
themselves} &
\textbf{This model} &
\textbf{Theorem~\ref{thm:entry}. Without it the model must give up either the
mechanism or a determinate prediction.}\\
\bottomrule
\end{longtable}
\end{center}

\subsection{Contributions}

Two things, stated precisely so that neither is oversold.

\begin{enumerate}
\item \textbf{A measurement result.} Work on ownership and tariffs has used
country-level ownership shares, because that is what the data allowed. With ownership
recorded firm by firm we show \emph{analytically} that the country-level share is not
merely imprecise but biased in a predictable direction, and we give the exact size of
the bias: it is the covariance between ownership and firm size
(Proposition~\ref{prop:main}). This is an identity, not an estimate.
\item \textbf{A uniqueness theorem.} Models with multi-factory firms and entry
normally have many equilibria, and the usual responses are to abandon the multi-factory
structure or to select an equilibrium by assumption. We do neither. We prove there is
exactly one, under a condition that is easy to state, easy to interpret, and directly
checkable in data (Theorem~\ref{thm:entry}).
\end{enumerate}

\newpage
\section{Six Stylized Facts}
\label{sec:factsstate}
\label{sec:facts}

These are referred to throughout by number, so they are stated here once. They come
from firm-level customs records matched to ownership data for ten Latin American
origin countries, 2006--2022. Section~\ref{sec:quant} compares each one with the model.

\begin{center}
\begin{tabular}{cp{11.6cm}}
\toprule
\# & The pattern in the data\\
\midrule
1 & International firms account for a large share of export value --- between 46\% and
74\% depending on the country --- and the great majority of that is \emph{foreign}
owned rather than locally owned.\\
\addlinespace
2 & Foreign-owned firms are concentrated in \emph{complex} products; locally-owned
international firms are concentrated in simple and primary goods.\\
\addlinespace
3 & The parent firms come from a small handful of countries.\\
\addlinespace
4 & Counting each factory as a separate firm understates how concentrated markets are.
Grouping factories by their ultimate owner raises measured concentration --- and export
value is dominated by a few very large groups.\\
\addlinespace
5 & In markets where more international firms are present, \emph{local} firms export
more, not less.\\
\addlinespace
6 & Distance discourages exports less for international firms than for others, and
least of all for a firm that already has a factory in the destination country.\\
\bottomrule
\end{tabular}
\end{center}

\noindent Facts 3, 4 and 6 are not used to choose any parameter, so the model's
agreement with them is a genuine prediction. Fact 5 is the one the model gets wrong,
and Section~\ref{sec:fact5} says so plainly.

\newpage
\section{Model}
\label{sec:model}

\subsection{Set-up}

There are $N$ countries and $K$ sectors. A \textbf{market} is one (destination,
sector) pair: one product sold in one country. Inside a market, several
\textbf{factories} compete. Each factory belongs to a \textbf{group} --- the firm that
ultimately owns it. A group may own one factory or many, and its factories may sit in
different countries.

Throughout, ``a firm's market share'' means the \emph{group's} share, adding up all
its factories in that market. That single choice is the heart of the model, and the
next section explains why.

\subsection{The mechanism: a group competes with itself}
\label{sec:cannibal}

Suppose a group owns two factories selling into the same market, and expands one of
them. Two things happen.

\begin{itemize}
\item It takes sales from rival firms. That is ordinary competition.
\item It \emph{also} takes sales from its own other factory. That loss is not
somebody else's problem: the group collects both sets of profits, so it counts the
loss against the gain.
\end{itemize}

\noindent We call the second effect \textbf{cannibalisation}. It makes a group hold
back and charge more than any single factory would. And the larger the group's share
already is, the more of any expansion it takes from itself. Hence:

\begin{center}
\emph{the larger a group's share, the higher the markup it charges.}
\end{center}

\begin{keypoint}
Big firms charge more because, when they grow, they are partly competing against
themselves. Every result in this document is a consequence of that one sentence. It
is also the reason the model cannot use the standard trade workhorse, in which all
firms charge the same markup no matter how large they are.
\end{keypoint}

\subsection{The markup}

\begin{align}
  \text{quantity competition:}\quad  \frac{1}{\mu(S)} &= 1 - \frac{1-S}{\sigma} - \frac{S}{\eta},
  \label{eq:muC}\\[4pt]
  \text{price competition:}\quad \mu(S) &= \frac{\varepsilon(S)}{\varepsilon(S)-1},
  \qquad \varepsilon(S)=\sigma-(\sigma-\eta)S .
  \label{eq:muB}
\end{align}

\says{$\mu$ is the markup: the ratio of price to marginal cost. A markup of $1.25$
means the price is 25\% above cost. Equation~\eqref{eq:muC} says the markup is set by
a weighted average of two elasticities. On the fraction $1-S$ of the market the group
does \emph{not} own, it faces the tough within-sector elasticity $\sigma$. On the
fraction $S$ it \emph{does} own, expanding only steals from itself, so it effectively
faces the much gentler across-sector elasticity $\eta$. As $S$ rises, more weight
falls on the gentle elasticity and the markup goes up. A negligible firm charges
$\sigma/(\sigma-1)$; a firm owning the entire market charges the monopoly markup
$\eta/(\eta-1)$. Equation~\eqref{eq:muB} is the same story when firms set prices
instead of quantities, and it runs between exactly the same two endpoints.}

\why{Because we need markups that \emph{vary with firm size}, and this is the standard
way to get them. Two elasticities are the minimum: a firm needs close rivals (other
suppliers of the same product) and distant ones (everything else the country buys),
and a single elasticity cannot express both. We write the model for a general markup
function $\mu(\cdot)$, with \eqref{eq:muC} and \eqref{eq:muB} as the two leading
cases, because as the box below explains, nothing important depends on which one is
used.}

\from{The two-elasticity structure and both markup formulas are Atkeson and Burstein
(2008). What is added here is applying $\mu$ to the \emph{group's} total share when a
group owns several factories in several countries --- which is what makes the
framework speak to multinational firms.}

\begin{keypoint}
\textbf{An objection worth answering before anyone raises it.} It is natural to read
\eqref{eq:muC} as ``the model assumes Cournot competition, so the results are an
artifact of that assumption''. They are not.

All the model requires is that a group \emph{internalises competition between its own
factories}. That is true under quantity competition and under price competition alike.
In both cases the group charges one markup over each factory's cost, and in both cases
that markup depends on the group's total share and nothing else. So every proof in
Section~\ref{part:proofs} is written for a general $\mu(\cdot)$ satisfying
Assumption~\ref{as:mu} below, and the program runs both. The books balance to the same
precision either way (Section~\ref{sec:accounting}), and the uniqueness condition is
in fact \emph{looser} under price competition.

What \emph{does} depend on the choice is the \emph{magnitude} of the final result. See
Section~\ref{sec:kappa}.
\end{keypoint}

\begin{assumption}[The markup rises with size]\label{as:mu}
$\mu$ is continuously differentiable and strictly increasing on $[0,1)$, with
$\mu(0)=\sigma/(\sigma-1)$ and $\sigma>\eta\ge1$.
\end{assumption}

\says{This is cannibalisation written formally: bigger groups charge more. If it
failed, large firms would charge \emph{less} than small ones, which is backwards, and
every proof in Section~\ref{part:proofs} would collapse. Both \eqref{eq:muC} and \eqref{eq:muB} satisfy
it whenever $\sigma>\eta$, which simply says goods within a sector are better
substitutes for each other than for goods outside it.}

\subsection{Profit}

\begin{equation}
  \Pi_g \;=\; \underbrace{E\,S_g}_{\text{sales}}\;\times\;\underbrace{L(S_g)}_{\text{margin}},
  \qquad L(S)\equiv 1-\frac1{\mu(S)} .
  \label{eq:profit}
\end{equation}

Under quantity competition the margin is exactly a straight line in the share,
$L(S)=(1+\kappa S)/\sigma$ with $\kappa\equiv\sigma/\eta-1$, so \eqref{eq:profit}
becomes
\begin{equation}
  \Pi_g=\frac{E}{\sigma}\,S_g\bigl(1+\kappa S_g\bigr).
  \label{eq:profitquad}
\end{equation}

\says{Profit is sales times the profit margin. The margin $L=1-1/\mu$ is the fraction
of the price that is not cost: a markup of $1.25$ means a margin of $0.20$. Because
the markup rises with size, so does the margin --- and therefore profit rises
\emph{faster than proportionally} with size. A group twice as large earns more than
twice as much, since it sells twice as much \emph{and} earns more on each unit.}

\why{The $S^2$ term in \eqref{eq:profitquad} is where the entire paper lives. It is
the reason \emph{which} firms are large, and who owns them, matters for a country's
income. If $\kappa$ were zero, profit would be a flat fraction of sales, ownership
would be a pure transfer, and the question in Section~1 would have no content ---
this is Proposition~\ref{prop:irrelevance}.}

\from{Equation~\eqref{eq:profit} is an accounting identity given the markup. The
linearity of $L$ under quantity competition is a property of \eqref{eq:muC} and is
what makes Proposition~\ref{prop:main} an exact identity rather than an
approximation.}

\subsubsection{\texorpdfstring{$\kappa$}{kappa}: the one number that sets the size of the result}
\label{sec:kappa}

\begin{center}
\begin{tabular}{lcc}
\toprule
 & quantity competition & price competition\\
\midrule
$\kappa$ & $\sigma/\eta-1$ & $(\sigma-\eta)/\sigma$\\
at $\sigma=5,\ \eta=1$ & $4.00$ & $0.80$\\
\bottomrule
\end{tabular}
\end{center}

\says{$\kappa$ measures how sharply profit rises with size. The entire ownership
correction in Section~\ref{sec:result} is proportional to it, so the choice of how firms compete
changes the \emph{size} of the headline number by a factor of five --- while changing
nothing structural.}

\why{We report quantity competition as the baseline for two checkable reasons, not
from preference. First, the margin is exactly linear there, which makes
Proposition~\ref{prop:main} an identity rather than a second-order approximation.
Second, the factory-level concentration it implies is closer to what the data show.
Price competition is fully supported and every headline is available under it, so a
reader sees a range rather than a single number chosen for convenience.}

\subsection{Market clearing}

\begin{equation}
  \sum_g S_g \;=\; 1 .
  \label{eq:clearing}
\end{equation}

\says{Everything sold in the market is sold by somebody. Trivial as an accounting
statement, but it is the equation that \emph{determines} the market outcome: given how
strong each group is, \eqref{eq:clearing} is one equation that pins down the single
number $A$ defined next, and every price, share and profit follows from $A$.}

\subsection{Strength, and the map from strength to share}

Combining Equation~1 with the definition of a market share gives, for \emph{any}
markup function,
\begin{equation}
  S_g=\mu(S_g)^{1-\sigma}\frac{K_g}{A},
  \qquad\text{where}\qquad
  K_g\equiv\sum_{i\in g}c_i^{1-\sigma},
  \qquad
  A\equiv\sum_j p_j^{1-\sigma}
  \quad\text{($j$ running over every \emph{factory} in the market).}
  \label{eq:share}
\end{equation}

Define the \textbf{share map} $\psi(S)\equiv S\,\mu(S)^{\sigma-1}$. Then
\eqref{eq:share} says $\psi(S_g)=K_g/A$, that is
\begin{equation}
  \boxed{\;S_g=\psi^{-1}\!\left(\frac{K_g}{A}\right).\;}
  \label{eq:psiinv}
\end{equation}

\says{$K_g$ --- we call it the group's \textbf{strength} --- adds up how cheaply each
of the group's factories can supply this market. A group with more factories, or
cheaper ones, has higher $K_g$. $A$ is a summary of how tough the market is: it adds
up all prices being charged, so a crowded market of cheap rivals has high $A$.
Equation~\eqref{eq:psiinv} then says something simple: \emph{a group's share depends
on its own strength relative to the toughness of the market, and on nothing else.}}

\begin{keypoint}
\textbf{$K_g$ and $A$ are the two objects that make this model workable.}

However many factories a group owns, and in however many countries, its entire
strategic position in a market is one number, $K_g$. Cannibalisation is \emph{already
inside} that number --- it is not an extra complication layered on top. And whatever
its rivals are doing, all the group needs to know about them is one number, $A$.

So a market containing hundreds of factories is really a problem in \emph{one}
unknown. This is why Section~\ref{part:proofs} can prove things that would otherwise be out of reach,
and why the program solves the whole world economy in minutes.
\end{keypoint}

\why{Writing the model in terms of $K_g$ rather than in terms of individual factories
is not a computational trick; it is what makes the entry problem tractable. In
Section~\ref{part:proofs} a group choosing which of $n$ factories to operate looks like a choice among
$2^n$ combinations. Because only $K_g$ matters, it collapses to a choice of a single
number, and Theorem~\ref{thm:entry} follows.}

\from{The collapse to a single index is a property of nested constant-elasticity
demand and is implicit in Atkeson and Burstein (2008). Naming it and using it as the
state variable for the \emph{entry} problem is what is new here.}

\subsection{Technology and delivered costs}

A factory $a$ whose group is headquartered in $h$, which produces in $\ell$ and sells
to $n$ in sector $k$, has delivered cost
\begin{equation}
  c_{a,n}=\frac{1}{\varphi_a}\,
  \underbrace{w_\ell^{\,1-\nu_k}}_{\text{(i) wages, host only}}\,
  \underbrace{\bigl(P^{\text{IO}}_{\ell k}\bigr)^{\nu_k}}_{\text{(ii) inputs}}\,
  \underbrace{\gamma_{h\ell}}_{\text{(iii) distance from HQ}}\,
  \underbrace{d_{\ell n}}_{\text{(iv) shipping}}\,
  \underbrace{(1+t_{\ell n k})}_{\text{(v) tariff}},
  \label{eq:cost}
\end{equation}
where $P^{\text{IO}}_{\ell k}=\prod_{k'}P_{\ell k'}^{\,\omega_{k'k}}$ is the price of
the bundle of inputs sector $k$ buys, at the prices prevailing in the country where
the factory sits. In the calibration $d_{\ell n}=\mathrm{dist}_{\ell n}^{\,1/(\sigma-1)}$,
which makes the model's gravity elasticity of \emph{trade} with respect to distance
exactly $-1$, the central estimate of the gravity literature --- so the model's own
gravity regression is a free diagnostic rather than a target (it returns $-0.93$).

\says{Read the five pieces in turn.
\textbf{(i)} Workers are hired where the factory is, at the host wage $w_\ell$.
\textbf{(ii)} A fraction $\nu_k$ of costs is goods bought from other firms, at the
prices prevailing where the factory is.
\textbf{(iii)} Producing away from head office is less efficient by a factor
$\gamma_{h\ell}$, which equals one at home. \emph{This} is where being multinational
changes costs.
\textbf{(iv)} Shipping is counted from where the good is \emph{made}, not from head
office.
\textbf{(v)} Then the tariff.
Finally $\varphi_a$ is how good the factory is; a higher $\varphi_a$ means lower cost,
and it is where the group's head office shows up --- see the next paragraph.}

\subsubsection{Where the head office is, and why it is not in the cost equation}
\label{sec:hqchoice}

\says{A group's head office is what makes its factories good. It is not a shipment of
services that each factory buys at its parent's wage. That distinction sounds like
hair-splitting and it is not: it decides whether the model has one equilibrium.}

Head-office services enter through $\varphi_a$, not through the cost bundle. The
sector's \emph{complexity} $\alpha_k$ --- head-office intensity, rising across
sectors --- does two things, both in the productivity draw:
\begin{itemize}
\item a stand-alone local firm, one with no parent group, carries a penalty
$e^{-\text{hq\_gap}\,\alpha_k}$, because head-office capability is exactly what it
cannot buy at arm's length;
\item a group's capability is scaled by $(1+\text{adv\_slope}\,\alpha_k)$, because
complex goods are where capability pays.
\end{itemize}

\why{The alternative --- $\alpha_k$ as the share of cost paid at the \emph{parent's}
wage, so that a subsidiary abroad literally pays $w_h^{\alpha_k}w_\ell^{1-\alpha_k}$
--- is the Antr\`as headquarter input, it is supported in the code
(\texttt{hq\_cost = true}), and every result is reported under both. It is not the
baseline, for three reasons, all of them measured rather than argued:
\begin{enumerate}
\item it is the \emph{sole} mechanism behind the counterexample to gross substitutes
in Section~\ref{sec:wagecounter}; switching it off buys Theorem~\ref{thm:simple};
\item it puts a wrong-sign force on Fact~2. A subsidiary pays its parent's \emph{high}
wage precisely in the sectors where $\alpha_k$ is large, so the wage channel pushes
foreign shares \emph{down} in complex goods --- against the fact. Removing it moves the
Fact~2 gradient from $+0.33$ to $+0.42$, against a target of $+0.43$;
\item with it removed the fitted multinational productivity edge is no longer needed at
all, so three internally calibrated parameters become two.
\end{enumerate}
A referee is entitled to ask which version is right. The honest answer is that both
are in the literature, that the choice matters for tractability and for one fact, and
that the paper shows both.}

\from{$\gamma_{h\ell}$ is Ramondo and Rodr\'iguez-Clare (2013), and is measured for the
car industry by Head and Mayer (2019). Counting $d_{\ell n}$ from the production
country --- so that one factory can serve many markets --- is Tintelnot (2017). The
input bundle with weights $\nu_k$ and $\omega_{k'k}$ follows Caliendo and Parro (2015).
$\alpha_k$ as the index of head-office intensity is Antr\`as (2003) and Antr\`as and
Helpman (2004), used here on the capability side rather than the cost side. The Pareto
capability draw and the size-dependent pool of potential parents are Gaubert and
Itskhoki (2021); groups choosing \emph{sets} of production locations is Yang (2023).}

\subsection{Entry costs}

\begin{equation}
  F_{a,n}=f_k\;\underbrace{w_\ell}_{\text{same labour mix as production}}
  \;\times\;\underbrace{f^{\mathrm{dist}}_{\ell n}}_{\text{distance multiplier}},
  \label{eq:fixed}
\end{equation}
where $f^{\mathrm{dist}}_{\ell n}\ge1$ optionally scales the cost of serving a
distant market and equals one everywhere in the baseline calibration.

\says{Selling into a market costs something fixed each period, regardless of how much
is sold --- registration, a distribution network, marketing. A factory operates in a
market only if the profit it earns there covers this cost. This is the equation that
makes \emph{which factories exist} an outcome of the model instead of an assumption.}

\why{The detail that matters is that $F$ is paid with the \emph{same mix of labour as
production}. This does three jobs at once. (1) It keeps the model unchanged if every
wage is multiplied by the same number, so treating one country's wage as the unit of
account is legitimate. (2) It makes entry costs \emph{real resources}, so they appear
as somebody's wage income. (3) It means the profits handed to owners are profits
\emph{after} the cost of entering. Drop any one of the three and the accounting in
Section~\ref{sec:accounting} stops adding up --- which is precisely why we check it
there.}

\from{A fixed cost of serving each market is Tintelnot (2017). The ranking of
potential entrants by cost, and the idea that the \emph{number} of possible firms
scales with country size, are Gaubert and Itskhoki (2021).}

\subsection{Ownership and income}

\begin{equation}
  X_n \;=\; \underbrace{w_nL_n}_{\text{wages}}
  \;+\;\underbrace{\sum_g \theta_{gn}\bigl(\Pi_g-\text{entry costs}_g\bigr)}_{\text{profits owned, net}}
  \;+\;\underbrace{T_n}_{\text{tariff revenue}} ,
  \qquad
  T_n=\sum_{k}\sum_{a}\frac{t_{\ell nk}}{1+t_{\ell nk}}\,
      \frac{r_{a,nk}}{\mu_{g(a)}},
  \label{eq:income}
\end{equation}
with $r_{a,nk}$ factory $a$'s sales in market $(n,k)$ --- duties are levied on customs
value, i.e.\ on sales \emph{deflated by the markup}, which is why high-markup suppliers
yield little revenue per unit of distortion. $T_n=0$ in the baseline (no tariffs are
imposed in the calibration; the object exists for the policy experiments).

\says{A country's income is what its workers earn, plus its share of the profits of
every group in the world \emph{after} those groups have paid their entry costs, plus
the tariffs it collects. $\theta_{gn}$ is the fraction of group $g$ owned by country
$n$; for each group these fractions add up to one across countries.}

\why{This is the equation the whole paper is about. Profits are earned where factories
sell, but they are \emph{received} where owners live, and $\theta$ is what separates
the two. In the baseline, $\theta$ assigns each group entirely to the country of its
ultimate parent, which is what firm-level ownership data identify. Note that profits
are net of entry costs: an owner receives what is left after the firm has paid to be
in the market.}

\from{Cross-border ownership entering national income is Brecher and Bhagwati (1981)
and the literature following it. What is different here is that $\theta$ varies
\emph{firm by firm} rather than being one number per country --- and Section~\ref{sec:result} shows
that this distinction changes the answer.}

\subsection{The pool of potential firms}
\label{sec:pool}

Firms are not written down by hand. Two rules generate them.

\begin{enumerate}
\item \textbf{Larger countries draw from a larger pool.} The number of \emph{possible}
groups headquartered in a country grows with the size of that country, at a rate
governed by $\zeta$. Crucially, \textbf{the distribution of ability they draw from is
identical in every country}. Nothing assumes rich countries breed better firms.
\item \textbf{Head-office capability cannot be purchased.} A stand-alone local firm
--- one with no parent group --- is at a disadvantage, and that disadvantage is larger
in sectors with a higher $\alpha_k$.
\end{enumerate}

\says{Rule 1 says rich countries end up with better firms for a purely statistical
reason: \emph{draw more times and your best draw is better}. Rule 2 says a firm
without a parent group is worse at exactly the things complex products require.}

\why{Both rules were adopted because the model failed without them, which is the only
respectable reason to adopt anything. With pools of equal size everywhere, the model
produced as many international groups headquartered in the developing countries as in
the rich ones --- the opposite of one of the sharpest patterns in the data
(Section~\ref{sec:fact1}). With no head-office disadvantage, the model got the pattern
across product types \emph{exactly backwards}, for a reason set out in
Section~\ref{sec:fact2} that is instructive rather than embarrassing.}

\from{Rule 1 is the granular-firms mechanism of Gaubert and Itskhoki (2021). Rule 2 is
the head-office-input logic of Antr\`as (2003) and Antr\`as and Helpman (2004),
applied on the cost side.}

\subsection{Equilibrium}

\begin{definition}[Equilibrium]\label{def:eq}
An equilibrium is a set of wages $\{w_n\}$, price indices $\{P_{nk}\}$, market sizes
$\{E_{nk}\}$, incomes $\{X_n\}$, a set of operating factories, and quantities, such
that
\begin{enumerate}[label=(\roman*)]
\item every group maximises its profit in every market, taking rivals' behaviour as
given;
\item shares add up to one in every market, equation~\eqref{eq:clearing};
\item no group wishes to open or close a factory, given equation~\eqref{eq:fixed};
\item price indices are consistent with the costs they generate, via
equation~\eqref{eq:cost};
\item every worker is employed in every country;
\item each country's budget balances, equation~\eqref{eq:income}.
\end{enumerate}
\end{definition}

\noindent Condition (v) for one country is implied by the others together with
budget balance, so one wage can be fixed as the unit of account. Section~\ref{part:proofs} shows that
(i)--(iv) each have exactly one solution, and that (v)--(vi) have a solution which is
verified to be unique.

\newpage
\section{Existence and Uniqueness of Equilibrium}
\label{part:proofs}

\begin{keypoint}
\textbf{Why this part exists.} A model that admits several outcomes predicts nothing.
If two different market structures are both equilibria, then ``the model says parent
firms come from a few countries'' is a statement about which equilibrium somebody
chose to report, not about the model. Two of the six facts this model is meant to
explain \emph{are} facts about market structure. So uniqueness here is not a
technicality: it is the difference between a result and a preference.

\smallskip
\textbf{The roadmap.} The model solves as a nested sequence of one-dimensional
problems, each with a single solution:

\smallskip
\begin{center}
\begin{tabular}{clcc}
\toprule
 & Layer & Exists? & Unique?\\
\midrule
1 & prices and sales inside one market & \textbf{proved} & \textbf{proved} (Thm.~\ref{thm:market})\\
2 & which factories operate & yes, up to whole numbers & \textbf{proved} (Thm.~\ref{thm:entry})\\
3 & input prices, given wages & \textbf{proved} & \textbf{proved} (Prop.~\ref{prop:prices})\\
4 & spending and incomes, given prices & \textbf{proved} & \textbf{proved} (Prop.~\ref{prop:income})\\
5 & wages & \textbf{proved} (Thm.~\ref{thm:exist}) & under checked conditions (Thm.~\ref{thm:contract}; see \S\ref{sec:notproved})\\
\bottomrule
\end{tabular}
\end{center}
\end{keypoint}

\subsection{One market}

\says{We show that a market with any number of firms has exactly one outcome. The
argument is short because of Equation~4: the whole market reduces to one number, $A$,
and one downward-sloping curve. A downward-sloping curve crosses a horizontal line
once.}

\begin{lemma}[Strength converts into share, one for one]\label{lem:psi}
Under Assumption~\ref{as:mu}, the share map $\psi(S)=S\,\mu(S)^{\sigma-1}$ is
continuous and strictly increasing on $[0,1)$, with $\psi(0)=0$. It therefore has a
continuous, strictly increasing inverse, and \eqref{eq:psiinv} determines $S_g$
unambiguously from $K_g/A$.
\end{lemma}

\begin{proof}
$\psi$ is a product of two factors. The first, $S$, is non-negative and strictly
increasing. The second, $\mu(S)^{\sigma-1}$, is strictly positive and strictly
increasing, because $\mu$ is strictly increasing (Assumption~\ref{as:mu}) and
$\sigma>1$. A product of two such factors is strictly increasing. Continuity and
$\psi(0)=0$ are immediate.
\end{proof}

\begin{lemma}[Each group has exactly one best move]\label{lem:br}
Fix the behaviour of every other group. Then group $g$'s profit is strictly concave in
its own production choices, so it has a unique best response, and that response
involves producing a strictly positive amount at each of its factories.
\end{lemma}

\begin{proof}
Let $\rho=(\sigma-1)/\sigma\in(0,1)$, let $X_g=\sum_{i\in g}q_i^{\rho}$ be an index of
the group's own output, and let $B_g$ be the same index for all rivals, taken as given.

\emph{(i) Own choices matter only through $X_g$.} Group revenue is
$R=D^{1/\eta}X_g(X_g+B_g)^{e}$ with $e=(\eta-\sigma)/((\sigma-1)\eta)$, where $D$ is
the demand shifter. This is cannibalisation as an algebraic identity: a group's own
products are perfect substitutes for one another inside $X_g$, so a choice with many
dimensions collapses to one number.

\emph{(ii) Revenue rises in $X_g$ at a diminishing rate.} Since $\sigma>\eta$ we have
$e<0$, and $e\ge-1$ with equality only when $\eta=1$. For $e\in(-1,0)$,
$\partial R/\partial X=(X+B)^{e-1}[(1+e)X+B]>0$ and
$\partial^2R/\partial X^2=e(X+B)^{e-2}[(1+e)X+2B]<0$. When $\eta=1$ the expression
becomes $R=DX/(X+B)$, with first derivative $DB/(X+B)^2>0$ and second derivative
$-2DB/(X+B)^3<0$. Either way, rising at a diminishing rate.

\emph{(iii) Cost rises in $X_g$ at an increasing rate.} Minimising
$\sum_{i\in g}c_iq_i$ subject to $\sum_{i\in g}q_i^{\rho}=X$ gives
$C(X)=X^{\sigma/(\sigma-1)}K_g^{-1/(\sigma-1)}$. The exponent exceeds one, so $C$ is
strictly convex. Note which constant appears: the group's strength $K_g$. That is not
a coincidence --- $K_g$ \emph{is} the group's cost-efficiency index, which is why the
same object governs both pricing and entry.

Profit is therefore (rising at a diminishing rate) minus (rising at an increasing
rate): strictly concave. So the first-order condition identifies the optimum rather
than merely being necessary, and there is only one. Finally, as any $q_i\to0$ revenue
falls like $q_i^{\rho}$ while cost falls like $q_i$; since $\rho<1$, producing a
little always beats producing nothing.
\end{proof}

\begin{theorem}[A market has exactly one equilibrium]\label{thm:market}
Take any set of operating factories with strictly positive costs, belonging to at
least two groups. Under Assumption~\ref{as:mu} the market has exactly one equilibrium.
\end{theorem}

\begin{proof}
Define $T(A)=\sum_g\psi^{-1}(K_g/A)$: the total market share implied by a candidate
value of $A$. By Lemma~\ref{lem:psi} each term falls as $A$ rises, so $T$ is
continuous and strictly decreasing.

When $A$ is very large, every ratio $K_g/A$ is near zero, so every share is near zero
and $T(A)<1$.

In the other direction, let $\bar\psi=\lim_{S\to1^-}\psi(S)$ --- infinite when
$\eta=1$, finite otherwise --- and set $A_1=\max_g K_g/\bar\psi$. As $A$ falls towards
$A_1$ the largest group's share approaches one. Since there are at least two groups
and every other group has a strictly positive share, $T$ has already exceeded one
before $A$ reaches $A_1$.

A continuous, strictly decreasing function that is above one somewhere and below one
elsewhere crosses one exactly once. That pins down $A$. Every other quantity ---
shares, markups, prices, sales, quantities --- follows from $A$ by a one-way formula,
so each is unique too. Existence is constructive, and Lemma~\ref{lem:br} confirms the
result is a genuine equilibrium in which no group wishes to deviate.
\end{proof}

\subsection{Entry: which factories operate}

\subsubsection{The problem}

Each group chooses which of its possible factories to run in a given market. Write a
choice as a set $T$ of factories, with
\[
  \underbrace{K(T)=\sum_{i\in T}c_i^{1-\sigma}}_{\text{strength it buys}},
  \qquad
  \underbrace{\Phi(T)=\sum_{i\in T}F_i}_{\text{entry costs it pays}} .
\]
Given the market index $A$, the group's payoff is
\begin{equation}
  V(T;A)=\Omega\bigl(K(T);A\bigr)-\Phi(T),
  \qquad
  \Omega(K;A)\equiv \Pi\bigl(\psi^{-1}(K/A)\bigr).
  \label{eq:omega}
\end{equation}

\says{$\Omega(K;A)$ is the profit a group earns if it buys strength $K$ in a market
whose toughness is $A$. The group then subtracts the entry costs of the factories it
switched on. Because of Equation~4, the group's decision --- nominally a choice among
$2^n$ combinations of factories --- is really a choice of a single number $K$.}

\why{Entry decisions of this kind normally have many possible outcomes, because each
firm's entry makes the market tougher for the others. Whether firm~A or firm~B ends up
in the market can be genuinely indeterminate. That is the problem this section solves,
and Section~\ref{sec:whyhard} explains why the two standard fixes are unavailable
here.}

\subsubsection{Why the two standard fixes do not work here}
\label{sec:whyhard}

\begin{itemize}
\item \textbf{Rank the possible entrants by cost and let the cheapest enter first.}
This delivers a unique outcome --- but only if every entrant is a \emph{separate}
firm. The moment a group owns several factories, switching on one of its \emph{own}
factories \emph{raises} its share instead of lowering it, the argument's key
requirement fails, and it collapses. Abandoning multi-factory groups means abandoning
cannibalisation, which is the mechanism the whole paper rests on.
\item \textbf{Fix the order in which firms move and report the resulting outcome.}
With $G$ groups there are $G!$ possible orders and no principle selects one. The
predicted market structure would then be a property of the chosen order rather than of
the model --- fatal when market structure is the object being explained.
\end{itemize}

\subsubsection{The resolution: separate the firm's own effect from its rivals'}

The requirement that fails asks whether an entrant's extra profit falls as the
\emph{total} number of firms rises. That question mixes two different things: what a
firm does to \emph{itself}, and what its \emph{rivals} do to it. Separated, each
behaves.

\begin{itemize}
\item[\textbf{(A)}] \textbf{Own effect:} a group's payoff rises with its own strength,
but at a diminishing rate.
\item[\textbf{(B)}] \textbf{Rival effect:} extra strength is worth \emph{less} when
the market is tougher.
\end{itemize}

\noindent The next two lemmas prove both. They are results, not assumptions. To state
them we need two elasticities:
\begin{equation}
  \varepsilon_G(S)\equiv\frac{S\,G'(S)}{G(S)}\ \ \text{with}\ \ G(S)\equiv\frac{\Pi'(S)}{\psi'(S)},
  \qquad
  \varepsilon_\psi(S)\equiv\frac{S\,\psi'(S)}{\psi(S)} .
  \label{eq:elast}
\end{equation}

\says{$G(S)$ is the extra profit a group gets from one more unit of strength. So
$\varepsilon_G$ measures whether that reward shrinks as the group grows --- condition
(A) is $\varepsilon_G<0$. And $\varepsilon_\psi$ measures how fast strength turns into
share. Condition (B) turns out to be a comparison between the two, as Lemma~\ref{lem:B}
shows.}

\begin{lemma}[Condition (A): diminishing returns to own strength]\label{lem:A}
$\Omega$ is strictly concave in $K$ if and only if $\varepsilon_G(S)<0$. Under
quantity competition with $\eta=1$ and $\sigma\ge2$, this holds at \emph{every} share
$S\in(0,1)$.
\end{lemma}

\begin{proof}
From \eqref{eq:psiinv}, $\partial S/\partial K=1/(A\psi'(S))$, so
$\Omega_K=\Pi'(S)/(A\psi'(S))=G(S)/A$ and $\Omega_{KK}=G'(S)/(A^2\psi'(S))$. Since
$\psi'>0$, the sign of $\Omega_{KK}$ is the sign of $G'$, hence of $\varepsilon_G$.

Now specialise. With $\eta=1$, \eqref{eq:muC} gives
$\mu(S)=\sigma/[(\sigma-1)(1-S)]$, so the margin is linear,
$L(S)=[1+(\sigma-1)S]/\sigma$, and setting $E=1$,
\[
  \Pi'(S)=\tfrac1\sigma\bigl[1+aS\bigr],\qquad
  \psi'(S)=C(1-S)^{-\sigma}\bigl[1+bS\bigr],
\]
where $a\equiv2(\sigma-1)$, $b\equiv\sigma-2$ and $C=(\sigma/(\sigma-1))^{\sigma-1}$.
Hence $\ln G=\text{constant}+\ln(1+aS)+\sigma\ln(1-S)-\ln(1+bS)$ and
\begin{equation}
  \varepsilon_G(S)=S\,f(S),\qquad
  f(S)\equiv\frac{a}{1+aS}-\frac{b}{1+bS}-\frac{\sigma}{1-S}.
  \label{eq:epsG}
\end{equation}
At $S=0$, $f(0)=a-b-\sigma=2\sigma-2-(\sigma-2)-\sigma=0$: it starts exactly at zero.
Differentiating,
\[
  f'(S)=-\frac{a^2}{(1+aS)^2}+\frac{b^2}{(1+bS)^2}-\frac{\sigma}{(1-S)^2}.
\]
For $\sigma\ge2$ we have $b\ge0$ and $a>b$, so $a/(1+aS)>b/(1+bS)\ge0$ and therefore
$a^2/(1+aS)^2>b^2/(1+bS)^2$, making $f'(S)<0$. A function that starts at zero and
strictly decreases is negative thereafter. So $f<0$, hence $\varepsilon_G=Sf(S)<0$.
\end{proof}

\begin{lemma}[Condition (B): strength is worth less in a tougher market]\label{lem:B}
$\Omega_{KA}<0$ if and only if $\varepsilon_G(S)+\varepsilon_\psi(S)>0$. Under quantity
competition with $\eta=1$,
\begin{equation}
  \varepsilon_G(S)+\varepsilon_\psi(S)
  =\frac{aS}{1+aS}-\frac{bS}{1+bS}+\frac{1-2S}{1-S},
  \label{eq:B}
\end{equation}
which is strictly positive at every share $S\le\tfrac12$, for every $\sigma>1$. At
exactly $S=\tfrac12$ its value is $1/\sigma$.
\end{lemma}

\begin{proof}
Differentiating $\psi(S)=K/A$ with respect to $A$ gives
$\partial S/\partial A=-\psi(S)/(A\psi'(S))$. Hence
\[
  \Omega_{KA}=\frac{\partial}{\partial A}\Bigl[\frac{G(S)}{A}\Bigr]
  =\frac{G'(S)}{A}\frac{\partial S}{\partial A}-\frac{G(S)}{A^2}
  =-\frac{1}{A^2}\Bigl[G'(S)\frac{\psi(S)}{\psi'(S)}+G(S)\Bigr].
\]
Since $\Pi'>0$ and $\psi'>0$ we have $G>0$; dividing the bracket by $G$ and
multiplying by $\varepsilon_\psi>0$ converts ``$\Omega_{KA}<0$'' into
``$\varepsilon_G+\varepsilon_\psi>0$''.

For the closed form, $\varepsilon_\psi=S\psi'/\psi=(1+bS)/(1-S)$. Adding this to
\eqref{eq:epsG} and using the identity $1+bS-\sigma S=1-2S$ gives \eqref{eq:B}.

Now the sign, which is elementary. The map $x\mapsto x/(1+x)$ is strictly increasing,
and $a>b$ because $2(\sigma-1)>\sigma-2$ for every $\sigma>0$. So the first two terms
of \eqref{eq:B} are together strictly positive. The third term, $(1-2S)/(1-S)$, is
non-negative precisely when $S\le\frac12$. A strictly positive number plus a
non-negative number is strictly positive.

At $S=\frac12$ the third term vanishes and the first two give
$\frac{a}{2+a}-\frac{b}{2+b}=\frac{\sigma-1}{\sigma}-\frac{\sigma-2}{\sigma}
=\frac1\sigma$.
\end{proof}

\begin{keypoint}
\textbf{Lemma~\ref{lem:B} is the crux, so here it is in plain words.} The condition
that delivers a unique outcome holds \textbf{whenever no single group controls more
than half of a market} --- and it holds with room to spare, because the margin at
exactly one half is $1/\sigma$, which at $\sigma=5$ means the condition is about 20\%
from binding.

This is not a technical assumption buried in an appendix. It is a plain statement that
\emph{can be checked in the data}: is any single corporate group above 50\% of a
product market? At the model's own solution the largest group share anywhere is
$0.344$ under quantity competition and $0.420$ under price competition.
\end{keypoint}

\begin{lemma}[Only the efficient choices matter]\label{lem:frontier}
$\Omega$ is strictly increasing in $K$. Hence a set $T$ can be optimal only if no
other set delivers at least as much strength for no more money. The candidates are
therefore the efficient sets, and they form a \emph{chain}: moving up it, strength and
cost rise together.
\end{lemma}

\begin{proof}
$\Omega_K=G(S)/A>0$; the remainder is the definition of an efficient set, and the
chain property holds because along it $K$ and $\Phi$ increase together.
\end{proof}

\says{A group never runs an expensive factory when a cheaper one gives it more
strength. So although it nominally faces $2^n$ combinations, only a handful are ever
candidates, and they line up on a single ladder from ``small and cheap'' to ``large
and expensive''. Its real decision is where to stop climbing --- a cutoff, but a
cutoff \emph{within its own list} rather than across firms.}

\begin{theorem}[The entry outcome is unique]\label{thm:entry}
Suppose Assumption~\ref{as:mu} holds; that no group's share exceeds one half in any
market \emph{on the comparison region} --- at every strength profile between any two
candidate outcomes, not only at the outcomes themselves; and that no group is exactly
indifferent between two rungs of its ladder (generic in the capability draws). Then
there is \emph{at most one} entry outcome. No order of moves is imposed, no outcome is
selected by hand, and groups continue to compete with themselves. (Existence is a
separate matter: entry is in whole factories, so an exact equilibrium need not exist;
the measured near-miss is reported in Section~\ref{sec:notproved}.)
\end{theorem}

\noindent\emph{The equilibrium concept, precisely.} An entry outcome is a strength
profile $\{K_g\}$ and an index $A$ such that each $K_g$ is optimal for group $g$
against $A$, and $A$ is the market index generated by all the choices together
(Step~2's map). \emph{On the one-half hypothesis:} the proof compares two candidate
outcomes and needs condition (B) at the points in between, which is why the theorem
asks for one half on the region rather than only at equilibria. The check has slack to
spare: the sufficient threshold is $\tfrac12$, the exact frontier at $\sigma=5$ is
$0.545$, and the largest share at the computed solution is $\approx0.30$ --- so shares
would have to nearly double along the comparison path before the hypothesis could
bind, and the brute-force verification below, which needs no such hypothesis at all,
finds one outcome in every one of 120 markets. The numerical verification scores every
deviation on the fully re-solved market, which is the stronger notion; the two agree
in every tested case.

\begin{proof}
\emph{Step 1: each group wants to be smaller when the market is tougher.} Take two
rungs of the ladder of Lemma~\ref{lem:frontier}, with $K_1<K_2$ and costs
$\Phi_1\le\Phi_2$, and define
\[
  D(A)\;=\;\bigl[\Omega(K_2;A)-\Omega(K_1;A)\bigr]-(\Phi_2-\Phi_1)
  \;=\;\int_{K_1}^{K_2}\Omega_K(K;A)\,dK-(\Phi_2-\Phi_1),
\]
the net gain from choosing the larger rung. By Lemma~\ref{lem:B}, $\Omega_{KA}<0$, so
$D'(A)<0$: the gain strictly falls as the market gets tougher. Hence the larger rung
is preferred when $A$ is low and, once abandoned, is never preferred again. Applying
this to every pair of rungs, the optimal choice $K^{*}_g(A)$ never increases as $A$
rises. (This is the standard single-crossing argument. The candidate set is a chain
because $\Omega$ is increasing in $K$ --- Lemma~\ref{lem:frontier}; Lemma~\ref{lem:A}
is what makes each group's problem one-dimensional in the first place.)

\emph{Step 2: the market gets tougher when any group gets stronger.} Given strengths
$\{K_g\}$, the market index $A(\{K_g\})$ solves $\sum_g\psi^{-1}(K_g/A)=1$. Raising one
$K_j$ while holding $A$ fixed pushes the left-hand side above one; the left-hand side
strictly decreases in $A$ (proof of Theorem~\ref{thm:market}); so restoring equality
requires a strictly higher $A$. Hence $A(\cdot)$ strictly increases in every argument.

\emph{Step 3: two outcomes are impossible.} Suppose $(A,\{K_g\})$ and
$(A',\{K'_g\})$ are both entry equilibria, with $A<A'$. Being an equilibrium means
$K_g=K^{*}_g(A)$, $K'_g=K^{*}_g(A')$, $A=A(\{K_g\})$ and $A'=A(\{K'_g\})$. By Step 1,
$K'_g\le K_g$ for every $g$. By Step 2, $A'=A(\{K'_g\})\le A(\{K_g\})=A$. This
contradicts $A<A'$. Hence no two distinct equilibria exist.
\end{proof}

\begin{corollary}[Large groups and small firms need only one theorem]\label{cor:tiers}
A single-factory local firm is the case in which a group's ladder has one rung. So
Theorem~\ref{thm:entry} covers the large international groups and the many small local
firms simultaneously, with the small firms making genuine decisions rather than
forming a passive backdrop.
\end{corollary}

\begin{keypoint}
\textbf{What this buys.} The model keeps cannibalisation --- the mechanism, and the
thing that makes Fact~4 work --- \emph{and} makes a determinate prediction about which
firms exist. The literature's two options were to keep one or the other.
\end{keypoint}

\checked{The closed forms \eqref{eq:epsG} and \eqref{eq:B} were checked against brute
numerical differentiation of $\Omega$ itself: they agree to $4\times10^{-5}$, the
accuracy of the numerical method, so the algebra is machine-verified rather than merely
typed. Condition (A) failed at none of the shares tested for $\sigma$ between 2 and 12;
condition (B) at none of the shares at or below one half; and the predicted margin of
$1/\sigma$ at exactly one half matched to the last digit.

Lemmas~\ref{lem:A} and \ref{lem:B} are proved for the case $\eta=1$ with quantity
competition. The ``no group above one half'' rule was therefore also checked
numerically for other values of $\eta$ and for price competition. It is sufficient in
every case, and the case proved above is the tightest one.

Independently, the conclusion of Theorem~\ref{thm:entry} was tested by brute force ---
enumerating \emph{every} combination, letting each group choose any set of its own
factories, and evaluating each deviation on a fully re-solved market. In 120 of 120
randomly drawn markets exactly one outcome survived, never more.}

\subsection{Input prices and incomes}

\begin{proposition}[Input prices]\label{prop:prices}
Given wages and a set of operating factories, input prices have exactly one solution,
and repeated recalculation converges to it geometrically.
\end{proposition}

\begin{proof}
Costs depend on prices and prices depend on costs, so this is a circular problem. It is
a \emph{contraction}: each round of recalculation shrinks the remaining error by a
fixed factor. Two facts establish this. First, a market's price index scales one for
one with the costs entering it: multiply every cost by $\lambda$ and all shares and
markups are unchanged, so every price, and the index, is multiplied by $\lambda$.
Second, in \eqref{eq:cost} log cost depends on log prices only through
$\nu_k\sum_{k'}\omega_{k'k}\log P_{\ell k'}$, with \emph{non-negative} weights
$\omega_{k'k}\ge0$ summing to one across sectors --- non-negativity is what turns the
weight sum into a bound on the sup-norm response (the same monotone-averaging fact as
the pass-through lemma). Together these bound the error contraction factor by
$\max_k\nu_k<1$. The contraction-mapping theorem then gives exactly one solution and
geometric convergence.
\end{proof}

\begin{proposition}[Spending and incomes]\label{prop:income}
Given prices and shares, spending and incomes are the unique solution of two linear
systems.
\end{proposition}

\begin{proof}
Spending satisfies $E=\varepsilon X+BE$, where $\varepsilon$ collects the shares of
final spending going to each market and $B$ collects, for each market, the value of
inputs it buys from every other market per unit of its own sales; income satisfies
\eqref{eq:income}. Both are
linear in $(E,X)$ once shares are known. Note that entry costs enter as a constant, not
as a new nonlinearity, which is why adding entry does not break linearity here. The
matrix $B$ has non-negative entries and each column sums to less than one, because a
market spends only the fraction $\nu_k$ of its \emph{cost} bill on inputs and cost is
itself a fraction $1/\mu<1$ of sales. A non-negative matrix whose columns sum to less
than one has spectral radius below one, so $(I-B)$ is invertible with a non-negative
inverse and the solution is unique. The income system follows the same argument, since
profits and tariffs are proper fractions of sales.
\end{proof}

\subsection{Wages}

\begin{theorem}[Equilibrium wages exist]\label{thm:exist}
Hold the set of operating factories fixed, and write $z(w)$ for the vector of
excess demand for labour --- demand minus the available workforce, country by
country. Then $z$ is continuous
on the interior of the wage simplex, unchanged when all wages are scaled by a common
factor, satisfies $w\cdot z(w)=0$, and grows without bound in any country whose wage
approaches zero. An equilibrium therefore exists.
\end{theorem}

\begin{proof}
\emph{Continuity} follows by composition: Theorem~\ref{thm:market} gives market
outcomes varying continuously with costs; Proposition~\ref{prop:prices} gives a
contraction whose fixed point varies continuously with its parameters;
Proposition~\ref{prop:income} is a linear solve with an invertible matrix.
\emph{Scale invariance}: multiplying all wages by a common factor multiplies all costs,
prices and incomes by that factor and leaves every share unchanged, so the
\emph{quantity} of labour demanded does not move. \emph{Adding up} is the accounting
identity of Section~\ref{sec:accounting}, verified below to hold away from equilibrium
as well as at it. \emph{Boundary behaviour}: as a country's wage approaches zero,
producing there becomes arbitrarily cheap relative to anywhere else, so demand for its
labour grows without bound. These are the standard conditions guaranteeing existence.
\end{proof}

\subsubsection{What uniqueness of wages would require}
\label{sec:wageunique}

\says{Existence is the easy half. To get uniqueness the usual route is to show that
labour in different countries is a \emph{substitute}: if one country's wage rises,
employers everywhere else should want \emph{more} workers, not fewer. If that is true,
two different wage vectors cannot both clear every labour market, and the argument is
three lines long. The whole difficulty is establishing it in a model where markups
move, and --- in the variant of Section~\ref{sec:hqchoice} where head-office services are
bought at the parent's wage --- where a foreign wage enters the cost of goods made at home.}

Write $W_n(w)=w_n\,\mathrm{LD}_n(w)$ for the total wage bill of country $n$, and
\[
  M_{nj}(w)\;=\;\frac{\partial \ln W_n}{\partial \ln w_j}.
\]
Because every price, cost and income in the model is homogeneous of degree one in
wages, each row of $M$ sums to one. Excess demand for labour is
$z_n=W_n/w_n-L_n$, so for $j\neq n$ we have
$\partial z_n/\partial\ln w_j=(W_n/w_n)\,M_{nj}$, and therefore
\begin{center}
\emph{gross substitutes} $\iff$ $M_{nj}>0$ for every $n\neq j$ $\iff$
\emph{a country's wage bill rises when any other country's wage rises.}
\end{center}

\begin{proposition}[Gross substitutes gives uniqueness]\label{prop:abh}
Suppose $M_{nj}(w)>0$ for all $n\neq j$ at every $w$ in a set $D$ that contains, with
any two of its points, every wage vector lying between them. Then $D$ contains at most
one equilibrium wage vector up to a common scaling.
\end{proposition}

\begin{proof}
Let $w$ and $w^{*}$ be equilibria in $D$ that are not proportional. Put
$\lambda=\max_j w^{*}_j/w_j$, attained at some country $n_0$, and replace $w^{*}$ by
$w^{*}/\lambda$; excess demand is unchanged by a common scaling, so this is still an
equilibrium, and now $w^{*}\le w$ with equality at $n_0$ and strict inequality
somewhere. Raise the wages from $w^{*}$ to $w$ one country at a time, leaving $n_0$
alone. Each step weakly raises $z_{n_0}$, and the step at a country where the
inequality was strict raises it strictly. Hence $z_{n_0}(w)>z_{n_0}(w^{*})=0$, which
contradicts $z_{n_0}(w)=0$.
\end{proof}

\noindent Only the region matters, not the whole positive orthant. Measure a wage
vector by its \emph{spread}, the log gap between its highest and lowest wage relative
to a fixed reference; the spread does not change when all wages are scaled. Every
vector on the path used in the proof has spread at most the sum of the spreads of its
two endpoints. So:

\begin{corollary}\label{cor:region}
If gross substitutes holds at every wage vector of spread at most $2R$, there is at
most one equilibrium of spread at most $R$.
\end{corollary}

\subsubsection{Lemma 3: a wage rise raises every price index, by less than one for one}

\says{This is the workhorse, and it is unconditional. Three things were supposed to
make this model's price indices behave badly: markups that move with market share,
firms buying inputs from each other in a loop, and factories paying a foreign
head office. Lemma~3 says none of them can make a price index fall when a wage rises.
The reason is simple once seen: at each step the response is an \emph{average} of the
responses one level down, and an average of numbers between zero and one is between
zero and one.}

\begin{lemma}[Pass-through]\label{lem:pass}
At any wage vector at which Layer~1 has a solution, the matrix
$\Lambda=\partial\ln P/\partial\ln w$, whose rows are indexed by (country, sector)
pairs, is non-negative and has row sums equal to one.
\end{lemma}

\begin{proof}
\emph{Step 1: costs.} From Equation~5,
\[
  \hat c_a \;=\; (1-\nu_k)\big[\alpha_k\hat w_h+(1-\alpha_k)\hat w_\ell\big]
             \;+\;\nu_k\sum_{k'}\omega_{k'k}\,\hat P_{\ell k'},
\]
where a hat denotes a proportional change. The coefficients are non-negative and add
to $(1-\nu_k)+\nu_k=1$. The head-office term sits inside this line: it changes
\emph{which} wage a factory is exposed to, not the fact that the weights are a
probability distribution.

\emph{Step 2: inside one market.} Let $\hat c_g=\sum_{a\in g}(s_a/S_g)\hat c_a$ be the
sales-weighted cost change of group $g$. Differentiating
$S_g=\mu(S_g)^{1-\sigma}K_g/A$ gives
\[
  \hat S_g \;=\; -\chi_g\big(\hat c_g-\hat P_{nk}\big),
  \qquad
  \chi_g=\frac{\sigma_k-1}{1+(\sigma_k-1)\varepsilon_\mu(S_g)}\;>\;0,
  \qquad \varepsilon_\mu=\frac{S\mu'(S)}{\mu(S)}\ \ge 0 ,
\]
the last inequality because markups rise with size. Shares add to one, so
$\sum_g S_g\hat S_g=0$, and therefore
\[
  \hat P_{nk} \;=\; \sum_g \xi_g\,\hat c_g,
  \qquad
  \xi_g=\frac{S_g\chi_g}{\sum_{g'}S_{g'}\chi_{g'}}\;>\;0,\qquad \sum_g\xi_g=1 .
\]
The price index moves with a \emph{weighted average} of the cost changes. This is
where the markup channel goes: it changes the weights $\xi_g$ and nothing else.

\emph{Step 3: stacking.} Substituting Step~1 into Step~2 gives
$\hat P=\Gamma\hat w+\mathcal N\hat P$ with $\Gamma\ge0$, $\mathcal N\ge0$ and
$(\Gamma+\mathcal N)\mathbf 1=\mathbf 1$; the row sums of $\mathcal N$ alone are
$\nu_k<1$. Hence $\rho(\mathcal N)<1$, $(I-\mathcal N)^{-1}=\sum_{t\ge0}\mathcal N^t\ge0$,
and
\[
  \Lambda=(I-\mathcal N)^{-1}\Gamma\;\ge\;0,
  \qquad
  \Lambda\mathbf 1=(I-\mathcal N)^{-1}(\mathbf 1-\nu)=\mathbf 1 . \qedhere
\]
\end{proof}

\begin{corollary}
Raising one country's wage raises every price index in the world, by a proportion
between zero and one; raising all wages together raises every price index in the same
proportion. In particular every factory's cost change $\hat c_a$ also lies between
zero and one.
\end{corollary}

\subsubsection{Lemma 4: what a factory's wage bill does when costs move}

\begin{lemma}[Decomposition]\label{lem:decomp}
Let $V_a$ be what factory $a$, owned by group $g$, spends on labour and head-office
services to serve market $(n,k)$. Then
\[
  \hat V_a \;=\;
  \underbrace{\hat E_{nk}}_{\text{the market grew}}
  \;-\;\underbrace{(\sigma_k-1)\,(\hat c_a-\hat c_g)}_{\text{reshuffling inside the group}}
  \;-\;\underbrace{\big[1-\varepsilon_\mu(S_g)\big]\chi_g\,(\hat c_g-\hat P_{nk})}_{\text{the group against its rivals}} .
\]
\end{lemma}

\begin{proof}
$\hat V_a=\hat E_{nk}+\hat s_a-\hat\mu_g$ with $\hat s_a=(1-\sigma_k)(\hat p_a-\hat
P_{nk})$ and $\hat p_a=\hat c_a+\hat\mu_g$, so $\hat V_a=\hat E_{nk}
-(\sigma_k-1)(\hat c_a-\hat P_{nk})-\sigma_k\hat\mu_g$. Now
$\hat\mu_g=\varepsilon_\mu\hat S_g$ and $\hat S_g=-\chi_g(\hat c_g-\hat P_{nk})$ from
Step~2 of Lemma~\ref{lem:pass}. Collecting terms and using
$\chi_g[1+(\sigma_k-1)\varepsilon_\mu]=\sigma_k-1$ gives the stated form.
\end{proof}

\begin{corollary}[The markup channel, and the same one half as Theorem~\ref{thm:entry}]
\label{cor:half}
The third term has the gross-substitutes sign --- a group whose costs rise by less than
the market's spends \emph{more} on factors --- if and only if
$\varepsilon_\mu(S_g)\le1$. Under quantity competition with $\eta=1$,
$\mu=\sigma/[(\sigma-1)(1-S)]$, so $\varepsilon_\mu=S/(1-S)$ and the condition is
exactly $S_g\le\frac12$. \textbf{One condition --- no group holds more than half of any
market --- carries both this theorem and Theorem~\ref{thm:entry}.}
\end{corollary}

\begin{corollary}[Head-office neutrality]\label{cor:hqneutral}
Suppose all of a group's factories report to the same head office --- true by
construction here, since the head-office country is a property of the parent. Then
that country is paid the fraction $\alpha_k$ of the group's \emph{entire} bill in the
market. Averaging the middle term of Lemma~\ref{lem:decomp}
over the group with weights $s_a/S_g$ gives zero, so
\[
  \hat V_g=\hat E_{nk}-\big[1-\varepsilon_\mu(S_g)\big]\chi_g(\hat c_g-\hat P_{nk}) .
\]
The reshuffling of production \emph{inside} a multinational --- the term the
head-office linkage was expected to break --- does not appear in the head-office
channel at all.
\end{corollary}

\noindent Corollaries~\ref{cor:half} and \ref{cor:hqneutral} dispose of the two
channels that had been blamed for the failure of the standard argument. What is left is
a race between the first term of Lemma~\ref{lem:decomp}, which is positive because a
wage rise anywhere makes the world's markets bigger in money terms, and the residual
exposure in the third. That race is the theorem's second hypothesis.

\subsubsection{Theorem 4: when the wage vector is unique}

Fix an ordered pair of countries $n\neq m$ and consider raising $w_m$ alone. Let
$\lambda_a=\hat c_a$, $\lambda_g=\hat c_g$ and $\Lambda_{n'k}=\hat P_{n'k}$ be the
responses delivered by Lemma~\ref{lem:pass}; all lie in $[0,1]$. Let $\omega_{a,n'}$ be
the share of country $n$'s wage bill that factory $a$ pays out of its sales in market
$(n',k)$, so that these shares add to one across all factories, all markets and both
the production and the entry-cost part. Define
\[
  \bar E^{(m)}_n=\sum \omega_{a,n'}\,\hat E_{n'k}\ (\text{entry-cost part: } \lambda^F_a),
  \qquad
  \Xi_{nm}=\sum \omega_{a,n'}\Big[(\sigma_k-1)(\lambda_a-\lambda_g)^{+}
            +\big[1-\varepsilon_\mu(S_g)\big]\chi_g(\lambda_g-\Lambda_{n'k})^{+}\Big],
\]
where $x^{+}=\max\{x,0\}$. In words: $\bar E^{(m)}_n$ is how much the markets country
$n$ earns in are growing, and $\Xi_{nm}$ is how much of country $n$'s income sits in
factories that are \emph{more} exposed to $w_m$ than their own group is, or in groups
more exposed than the market they sell in.

\begin{assumption}[No group above one half]\label{ass:H1}
$\varepsilon_\mu(S_g)\le1$ in every market. Under quantity competition with $\eta=1$
this is exactly: no parent holds more than half of any market.
\end{assumption}

\begin{assumption}[Demand beats exposure]\label{ass:H2}
$\bar E^{(m)}_n>\Xi_{nm}$ for every ordered pair $n\neq m$.
\end{assumption}

\begin{theorem}[Equilibrium wages are unique]\label{thm:wage}
Hold the set of operating factories fixed and let $\eta=1$. If
Assumptions~\ref{ass:H1} and \ref{ass:H2} hold at every wage vector of spread at most
$2R$, then gross substitutes holds there, and there is at most one equilibrium wage
vector of spread at most $R$.
\end{theorem}

\begin{proof}
Raise $w_m$ alone. By Lemma~\ref{lem:pass} every $\lambda$ and every $\Lambda$ lies in
$[0,1]$. Summing Lemma~\ref{lem:decomp} over the factories that pay country $n$, with
the weights $\omega_{a,n'}$,
\[
  M_{nm}=\sum\omega_{a,n'}\Big[\hat E_{n'k}-(\sigma_k-1)(\lambda_a-\lambda_g)
        -\big(1-\varepsilon_\mu\big)\chi_g(\lambda_g-\Lambda_{n'k})\Big]
        \;+\;\text{(entry-cost part)} .
\]
Since $\varepsilon_\mu\ge0$ we have $\chi_g\le\sigma_k-1$, and
Assumption~\ref{ass:H1} puts $1-\varepsilon_\mu$ in $[0,1]$; hence both reallocation
coefficients are non-negative, and each negative contribution is bounded below by minus
that coefficient times the positive part of its own gap. Therefore
$M_{nm}\ge\bar E^{(m)}_n-\Xi_{nm}>0$ by Assumption~\ref{ass:H2}. That is gross
substitutes; Proposition~\ref{prop:abh} and Corollary~\ref{cor:region} finish the
argument.
\end{proof}

\checked{Every analytic step above is checked against numerical differentiation of the
solved model by \texttt{verify\_wage\_theorem}, exactly as
\texttt{verify\_conditions\_analytic} does for Theorem~\ref{thm:entry}. At the
calibrated economy the pass-through matrix of Lemma~\ref{lem:pass} reproduces the
numerical derivative to $1.2\times10^{-7}$ with row sums equal to one to
$2.2\times10^{-15}$, the share response of Lemma~\ref{lem:pass} step~2 to
$2.2\times10^{-7}$, the three-term decomposition of Lemma~\ref{lem:decomp} to
$2.1\times10^{-7}$, and the implied country wage bill to $1.4\times10^{-7}$ --- the
error of the difference quotient itself.}

\subsubsection{The second hypothesis cannot be dropped}
\label{sec:wagecounter}

\says{An honest theorem says what it rules out. Assumption~\ref{ass:H2} is not a
technical nicety: without it the conclusion is false, and the way it fails is
economically specific. It fails when a country earns its living almost entirely from
head-office payments for factories located in one other country.}

\begin{proposition}[Gross substitutes is not a theorem of this model]
\label{prop:wagecounter}
There is a three-country economy satisfying every assumption of the model, and
satisfying Assumption~\ref{ass:H1} with room to spare, in which
$M_{23}<0$ at the equilibrium: country 2's wage bill \emph{falls} when country 3's
wage rises.
\end{proposition}

\noindent The economy: one sector; country 1 has local firms; country 2 owns parent
firms whose only factories are in country 3; country 3 supplies the workers for those
factories. Every euro country 2 earns is a head-office payment for production abroad.
Raise $w_3$: those parents become dearer, lose share to country 1's local firms, and
country 2's entire income falls with them. Each version of the example carries enough
local rivals that no parent holds more than a third of any market --- the largest
parent share ranges from $0.18$ to $0.33$ --- so Assumption~\ref{ass:H1} is satisfied
throughout and cannot be what fails. It is Assumption~\ref{ass:H2}, and it fails badly:
the smallest off-diagonal of $M$ reaches $-1.43$.

The statistic that separates the counterexample from the calibrated world is the
\textbf{bilateral exposure} $\psi_{nm}$: the share of country $n$'s labour income paid
by factories tied to country $m$, as host or as parent. In the counterexample
$\psi_{23}=1$. In the calibrated economy the largest bilateral exposure anywhere is
$0.38$ under quantity competition and $0.49$ under price competition. This is the
sharp form of the limit, and it is worth stating precisely, because the obvious guess
is wrong: \textbf{it is not that $\alpha_k$ is too large.} Gross substitutes survives
$\alpha_k=0.9$ in ordinary ownership geographies, and inside the counterexample family the
failure is \emph{worst} at $\alpha_k=0.15$ and has disappeared by $\alpha_k=0.80$. A large
head-office share is not the problem. The problem is a country whose entire income is a
thin slice of firms whose costs are set by one other country --- income \emph{concentrated}
on a single foreign partner.

\subsubsection{A version of the model in which the second hypothesis is not needed}
\label{sec:simple}

\says{Assumption~\ref{ass:H2} is the one with a counterexample, and the counterexample
runs on a single modelling choice: that a factory pays its \emph{parent's} wage on the
head-office share of its costs. That choice is doing two different jobs here, and only
one of them is needed. Separate them and the obstacle disappears.}

In the model as built, $\alpha_k$ appears twice. It is the share of a factory's cost
paid at its parent's wage --- the head-office input --- and it is the index of
\emph{complexity} in the productivity draws, through the penalty $e^{-\text{hq\_gap}\,
\alpha_k}$ on stand-alone local firms and the capability gradient
$\xi(1+\text{adv\_slope}\,\alpha_k)$ on parents. Only the first creates a cross-country
factor linkage, and only the first appears in the counterexample. The second is what
Facts~\ref{sec:fact1} and \ref{sec:fact2} are built on.

\begin{definition}[Head-office services as a capability]
Head-office services are a non-rival capability of the group rather than a factor
bought at the parent's wage: $\alpha_k=0$ in the cost function, while the productivity
draws keep the same $\alpha_k$ gradient. In the code this is
\texttt{world\_economy(...; hq\_cost = false)}; the productivity draws are bit-for-bit
identical to the baseline, so nothing behind Facts~1 and 2 is disturbed.
\end{definition}

\begin{theorem}[Gross substitutes without the second hypothesis]\label{thm:simple}
Suppose head-office services are a capability, intermediate inputs are absent
($\nu_k=0$), $\eta=1$, and Assumption~\ref{ass:H1} holds. Then for every ordered pair
$n\neq m$, every factory paying country $n$ satisfies
\[
  -(\sigma_k-1)(\lambda_a-\lambda_g)-\big[1-\varepsilon_\mu(S_g)\big]\chi_g(\lambda_g
  -\Lambda_{n'k})\;\ge\;(\sigma_k-1)\min\{\lambda_g,\Lambda_{n'k}\}\;\ge\;0 ,
\]
so the entire reallocation block is non-negative and
$\partial\ln W_n/\partial\ln w_m\ \ge\ \bar E^{(m)}_n$. \emph{Strict} gross
substitutes, and hence uniqueness by Proposition~\ref{prop:abh}, follows whenever the
demand term is strictly positive --- some market country $n$ earns in strictly grows
($\bar E^{(m)}_n>0$; if no market's spending falls the conclusion holds weakly). At
the calibrated economy the demand term clears zero with margin $+0.024$.
\end{theorem}

\begin{proof}
With $\alpha_k=0$ a factory's entire factor bill is paid to its host, so every factory
paying $n$ is located in $n$; with $\nu_k=0$ its cost responds only to its host's wage,
and $n\neq m$, so $\lambda_a=0$. Assumption~\ref{ass:H1} gives
$0\le[1-\varepsilon_\mu]\chi_g\le\sigma_k-1$, the upper bound because
$\varepsilon_\mu\ge0$. Writing $c=[1-\varepsilon_\mu]\chi_g$, the displayed expression is
$(\sigma_k-1)\lambda_g-c(\lambda_g-\Lambda)$. If $\lambda_g\ge\Lambda$ this is at least
$(\sigma_k-1)\lambda_g-(\sigma_k-1)(\lambda_g-\Lambda)=(\sigma_k-1)\Lambda$; if
$\lambda_g<\Lambda$ both terms are non-negative. Both cases are bounded below by
$(\sigma_k-1)\min\{\lambda_g,\Lambda\}$.
\end{proof}

\begin{remark}[What dropping input--output linkages buys, exactly]
With $\nu_k>0$ the same conclusion holds under one extra condition,
$\lambda_a\le\min\{\lambda_g,\Lambda_{n'k}\}$: country $n$'s input prices do not rise by
more than its group's average or its market's. Setting $\nu_k=0$ makes that condition
vacuous --- $\lambda_a$ is then identically zero. So the input--output block is not
what breaks anything, but it is what stands between a conditional statement and an
unconditional one. It is also the block added for Fact~\ref{sec:fact5}, which the model
fails with or without it, so switching it off costs nothing that currently works.
At the calibrated economy the floor moves from $-1.79$ (head-office cost,
input--output on) to $-0.11$ (capability, input--output on) to $+0.014$
(capability, input--output off), and the violation of the extra condition falls
from $0.45$ to $0.030$ to exactly zero.

\textbf{At $\nu=0$ the theorem certifies the calibrated economy.} Both of its inputs
hold on a box, not just the first: Assumption~\ref{ass:H1} with slack $0.54$ and the
demand condition $\bar E^{(m)}_n\ge0$ with slack $+0.024$, at every wage vector of spread
up to $0.80$. Gross substitutes follows there, and Corollary~\ref{cor:region} gives
\textbf{at most one equilibrium of spread $\le0.40$}. The baseline keeps $\nu=0.55$
because intermediates are $55$\% of gross output and the policy question is how a tariff
propagates; there the same margin is verified rather than certified --- gross substitutes
holds with margin $+0.147$ and never failed at any point tested. The facts barely move
between the two: Fact~1 goes $0.689\to0.727$ and Fact~2's gradient $+0.42\to+0.47$
against a target of $+0.43$.
\end{remark}

\begin{remark}[What is still not proved, and it is now the textbook one]
Theorem~\ref{thm:simple} removes the channel peculiar to this model. What remains is the
ordinary income effect: a country whose firms lose profits when a foreign wage rises
spends less, and the markets it sells into shrink. That is the classical obstacle to
gross substitutes in any model with non-labour income, it is not specific to
multinationals, and it is what the last column of the stress table measures.
\end{remark}

\checked{At the calibrated economy the smallest off-diagonal of $M$ --- the
margin by which gross substitutes holds --- widens from $+0.100$ under the baseline to
$+0.152$ when head-office services become a capability and $+0.360$ when the
input--output block is switched off as well, and gross substitutes survives out to a
spread of $2.0$ in all three. On deliberately adversarial random economies --- scrambled
cross-ownership, countries with no domestic production, extreme size asymmetry, tariffs
--- restricted to points where Assumption~\ref{ass:H1} holds, the baseline violates
gross substitutes at $19$ of $250$ points while the simplified model violates it at
$0$ of $340$. The exception is $\sigma=12$, far outside the calibrated range, where the
income effect bites at $3$ of $36$: Assumption~\ref{ass:H1} alone is not sufficient
in general, and that is what the previous remark says.}

\subsubsection{Uniqueness the constructive way: the model's own tâtonnement contracts}
\label{sec:contraction}

\says{Everything so far rules out a second equilibrium by importing a condition from
outside the model --- gross substitutes, then a classical argument. This subsection does
something more direct and, we think, more convincing: it shows that the map the
\emph{solver iterates} is a contraction. The model then has one solution because of how
it is solved, and the algorithm is guaranteed to find it from any starting point. The
only thing imported is Birkhoff's (1957) theorem, which is exact.}

Given a wage vector, everything else in the model is already a unique function of it:
market outcomes by Theorem~\ref{thm:market}, the operating set by
Theorem~\ref{thm:entry}, input prices by Proposition~\ref{prop:prices}, spending and
incomes by Proposition~\ref{prop:income}. So the whole model collapses to one map on
wages. The solver iterates
\begin{equation}
  \Psi_\kappa(w)_n \;=\; w_n\Bigl(\frac{\mathrm{LD}_n(w)}{L_n}\Bigr)^{\!\kappa}
  \;=\; w_n^{\,1-\kappa}\Bigl(\frac{W_n(w)}{L_n}\Bigr)^{\!\kappa},
  \qquad \kappa\in(0,1),
  \label{eq:tat}
\end{equation}
raising a country's wage when its labour is in excess demand. Its fixed points, up to a
common scaling, are exactly the equilibrium wage vectors.

$\Psi_\kappa$ is homogeneous of degree one, so it acts on the \emph{projective} space of
wage vectors --- wage vectors that differ only by a common factor are the same point.
The natural distance there is Hilbert's:
\[
  d(w,w')\;=\;\max_n\ln\frac{w_n}{w'_n}\;-\;\min_n\ln\frac{w_n}{w'_n},
\]
which is exactly the ``spread'' used earlier and is unchanged if either vector is
rescaled. The log-Jacobian of $\Psi_\kappa$ is
\[
  A_\kappa \;=\; (1-\kappa)I+\kappa M,\qquad M_{nj}=\frac{\partial\ln W_n}{\partial\ln w_j},
\]
and $M$ has \emph{unit row sums} --- that is the model's own homogeneity, not an
assumption --- so $A_\kappa$ does too.

\begin{theorem}[The tâtonnement is a contraction]\label{thm:contract}
Let $D$ be a set of wage vectors that is \emph{convex in logarithms} (the boxes used
throughout are), and suppose that on $D$
\begin{enumerate}
\item[(i)] $M_{nj}>0$ for all $n\neq j$ \quad(gross substitutes), and
\item[(ii)] $\displaystyle \kappa<\frac{1}{1-\min_{w\in D}\min_n M_{nn}(w)}$ \quad(enough damping).
\end{enumerate}
Then $A_\kappa(w)$ is a strictly positive stochastic matrix for every $w\in D$, and
for any $w,w'\in D$,
\[
  d\bigl(\Psi_\kappa(w),\Psi_\kappa(w')\bigr)\;\le\;\kappa_B\,d(w,w'),
  \qquad
  \kappa_B\;=\;\sup_{v\in D}\,\tanh\!\bigl(\Delta(A_\kappa(v))/4\bigr)\;<\;1,
  \qquad
  \Delta(A)=\max_{i,j,k,l}\ln\frac{A_{ik}A_{jl}}{A_{jk}A_{il}} .
\]
Consequently $D$ contains \textbf{at most one} equilibrium wage vector up to scaling;
and if in addition $\Psi_\kappa$ maps $D$ into itself, the iteration converges to that
equilibrium geometrically at rate $\kappa_B$ from any start in $D$.
\end{theorem}

\begin{proof}
$W$ is homogeneous of degree one in $w$, so $\Psi_\kappa$ is too and the rows of $M$,
hence of $A_\kappa$, sum to one. Off the diagonal $A_{\kappa,nj}=\kappa M_{nj}>0$ by (i);
on the diagonal $A_{\kappa,nn}=(1-\kappa)+\kappa M_{nn}>0$ exactly when
$\kappa(1-M_{nn})<1$, which is (ii). Work in log coordinates, where Hilbert's metric is
the oscillation seminorm $\mathrm{osc}(x)=\max_n x_n-\min_n x_n$ and $\Psi_\kappa$
becomes a map $\varphi$ with Jacobian $A_\kappa$. Three steps.
\emph{First}, by the integral form of the mean value theorem,
$\varphi(x)-\varphi(x')=\bar A\,(x-x')$ with $\bar A=\int_0^1 A_\kappa(x'+t(x-x'))\,dt$
--- and the segment stays in $D$ because $D$ is log-convex, so $\bar A$ is an average
of strictly positive stochastic matrices, hence itself one.
\emph{Second}, the oscillation contraction of a stochastic matrix is measured by
Dobrushin's coefficient $\delta(A)=\tfrac12\max_{i,j}\lVert A_{i\cdot}-A_{j\cdot}\rVert_1$,
which is \emph{convex} in the matrix (a maximum of norms of linear functions), so
$\delta(\bar A)\le\int_0^1\delta\bigl(A_\kappa(t)\bigr)dt\le\sup_{v\in D}\delta(A_\kappa(v))$.
\emph{Third}, pointwise, $\delta(A)\le\tanh(\Delta(A)/4)$ --- the additive shadow of
Birkhoff's (1957) projective bound. Chaining the three,
$\mathrm{osc}(\varphi(x)-\varphi(x'))\le\kappa_B\,\mathrm{osc}(x-x')$. Two fixed points
in $D$ would contradict this directly, which gives uniqueness in $D$ with no further
hypothesis; convergence of the iteration additionally needs the iterates to remain in
$D$, whence the invariance clause. (In practice the solver's iterates are observed to
stay well inside the certified boxes, and ten dispersed starts reach the same wages;
invariance is a property checked on the run, not assumed.)
\end{proof}

\says{In ordinary words: measure the gap between two wage vectors by the spread of
their log-ratio. One application of the solver's update shrinks that gap by a definite
factor whenever, everywhere between the two vectors, a rise in any one country's wage
raises every \emph{other} country's wage bill (gross substitutes) and the update step
is damped enough. Two different equilibria would have to keep their distance under a
map that strictly shrinks all distances --- impossible --- so the region holds at most
one. The three-step argument in the proof is only the careful version of ``the average
of shrinking maps shrinks''.}

\begin{remark}[The damping is not a numerical convenience]
Condition (ii) has real content. Labour demand is \emph{elastic}: at the calibrated
economy $M_{nn}\approx-1.49$, so a country's own wage bill falls by about $1.5\%$ when
its wage rises $1\%$, and the \emph{undamped} map $\kappa=1$ is not order-preserving at
all. The model itself says how much damping is needed --- here
$\kappa<1/(1+1.49)=0.402$ --- and the solver's $\kappa=0.25$ sits inside it. With the
input--output block switched off the requirement tightens to $\kappa<0.272$, so $0.25$
is still admissible but with much less room; $0.15$ would be the safer default there.
\end{remark}

\checked{At the calibrated five-country economy (the input--output row is the same
economy with $\nu=0$), scanning boxes of wage vectors of increasing spread:

\begin{center}
\begin{tabular}{lrrrrrr}
\toprule
 & spread & $\min M_{nn}$ & largest $\kappa$ & $\min A_\kappa$ & $\Delta$ & $\kappa_B$\\
\midrule
quantity competition & 0.0 & $-1.489$ & 0.402 & $+0.0367$ & 4.65 & \textbf{0.822}\\
                     & 0.8 & $-1.669$ & 0.375 & $+0.0296$ & 5.57 & 0.884\\
                     & 2.0 & $-1.667$ & 0.375 & $+0.0075$ & 7.14 & 0.945\\
price competition    & 0.0 & $-1.568$ & 0.390 & $+0.0350$ & 4.65 & 0.822\\
                     & 2.0 & $-1.760$ & 0.362 & $+0.0077$ & 7.02 & 0.942\\
no input--output     & 0.0 & $-2.669$ & 0.273 & $+0.0827$ & 1.02 & \textbf{0.249}\\
\bottomrule
\end{tabular}
\end{center}

\noindent The map contracts on every box tested, under both ways of competing, out to a
spread of $2.0$ --- a factor of $e^{2}$ between the highest and lowest relative wage.
And the conclusion is checked rather than trusted: taking random \emph{pairs} of wage
vectors and applying the solver's map to both, the Hilbert distance shrank by a factor
of at most $0.429$, comfortably inside the certified bound of $0.872$, in every one of
the pairs tried.}

\subsubsection{How far the theorem reaches, and what is still only checked}

Two things are still evidence rather than proof, and both are stated rather than
buried.

\begin{enumerate}
\item \textbf{The hypotheses are verified on a region, not everywhere.}
Assumptions~\ref{ass:H1} and~\ref{ass:H2} are inequalities in objects the solver
computes, so they are checked, not assumed --- the same standing as
Theorem~\ref{thm:entry}'s condition (B). At the calibrated equilibrium
Assumption~\ref{ass:H1} holds with the largest parent share at $0.34$ against a
threshold of $0.50$, and Assumption~\ref{ass:H2} holds for all $20$ ordered pairs
under both ways of competing. Away from the equilibrium they are checked on a grid of
the surrounding region --- and, in the $\nu=0$ version, the gap between grid points is
now closed: an outward-rounded interval-arithmetic computation (Moore 1966) encloses the
gross-substitutes matrix over the \emph{whole} of a box of wage vectors around the
solution at once. The certified region reaches a spread of $0.30$ --- every wage
vector within $\pm15\%$ (log) of the solved relative wages, enclosed as a continuum
--- so there ``exactly one equilibrium'' is a theorem with no between-points gap.
Beyond that box the grid certificate (out to spread $2.0$) is what remains.

\item \textbf{Gross substitutes is sufficient, not necessary.} What uniqueness actually
needs is that every equilibrium have index $+1$ --- the sign of the determinant of the
reduced Jacobian equal to $(-1)^{N-1}$ --- since the indices of all equilibria must sum
to one. Walras' law makes the vector of wage bills the left null vector of $I-M$, so
every diagonal cofactor of $I-M$ is proportional to a wage bill and \emph{the index does
not depend on which country is chosen as the numeraire}. In $50$ random economies drawn
from the family of Proposition~\ref{prop:wagecounter}, gross substitutes fails at the
equilibrium in $34$ --- and the index is $+1$ in all $50$, is numeraire-independent in
all $50$, and $25$ dispersed starting points find a single equilibrium in every one of
them. So the counterexample kills the standard \emph{argument}, not the
\emph{conclusion}.
\end{enumerate}

\subsection{The model really is in general equilibrium}
\label{sec:accounting}

\says{``General equilibrium'' means nothing leaks. Every euro of sales is somebody's
wage, somebody's input purchase, somebody's tariff receipt or somebody's profit; every
euro of income is spent; every worker has a job. It is easy to write a model that looks
like this and quietly loses money somewhere. So we check the books \emph{away from}
equilibrium, where an error has nowhere to hide, and under both ways of competing.}

\begin{center}
\begin{tabular}{lrr}
\toprule
Worst error over three very different wage vectors & quantities & prices\\
\midrule
Sales $=$ wages $+$ inputs $+$ tariffs $+$ profit & $3.0\times10^{-16}$ & $9.0\times10^{-16}$\\
Total wage bill $=$ production wages $+$ wages to enter markets & $1.2\times10^{-15}$ & $2.4\times10^{-15}$\\
\textbf{Adding-up condition, away from equilibrium} & $4.8\times10^{-16}$ & $5.9\times10^{-16}$\\
World income $=$ world spending & $9.5\times10^{-16}$ & $3.9\times10^{-16}$\\
Each country: spending $=$ wages $+$ profits owned, net $+$ tariffs & $2.6\times10^{-16}$ & $4.0\times10^{-16}$\\
\midrule
Same five checks at the solved equilibrium & $\le1.0\times10^{-15}$ & $\le1.1\times10^{-15}$\\
\bottomrule
\end{tabular}
\end{center}

\noindent The second line is a genuine test rather than a formality. Entry costs are
real resources, so they must appear \emph{both} as somebody's wage income \emph{and} as
a deduction from the profits handed to owners. Omit either half and the adding-up
condition fails.

\subsection{What is \emph{not} proved, and what was done instead}
\label{sec:notproved}

Everything that is checked numerically rather than proved is collected here --- the
two headline items below, plus the smaller ones flagged in place (the entry lemmas
for $\eta\neq1$ and price competition, and the away-from-equilibrium grids) --- so
none of it is left for a reader to find.

\subsubsection{That equilibrium wages are unique \emph{without conditions}}

That they exist is proved (Theorem~\ref{thm:exist}). That there is only one set of them
is proved three ways, all of them conditional on something that is \emph{checked} rather
than assumed, and the strongest is the constructive one:

\begin{itemize}
\item Theorem~\ref{thm:contract} --- the map the solver iterates is a contraction in
Hilbert's metric, so it has one fixed point and the algorithm converges to it. This is
the statement to quote: uniqueness follows from how the model is solved.  It needs the
tâtonnement Jacobian to be entrywise positive, which is gross substitutes plus a bound on
the damping that the model itself supplies.
\item Theorem~\ref{thm:simple} --- in the baseline, one condition (no group above half
a market) signs the whole reallocation block, leaving only a demand condition.
\item Theorem~\ref{thm:wage} --- the general statement, whose second hypothesis has a
counterexample and is the reason the baseline is what it is.
\end{itemize}

\noindent What is \emph{not} available is an unconditional theorem, and the reason is
worth stating plainly: the last channel that can break the conditions is the ordinary
income effect --- a country whose firms lose profits when a foreign wage rises spends
less --- and that channel is the subject of this paper. Removing it would mean spending
profits where they are earned rather than where they are owned, which deletes the
contribution. So the chain of simplifications stops here on purpose.

\says{The usual route to uniqueness is to show that goods are gross substitutes:
raising one country's wage should raise demand for everyone else's labour. Two things
were expected to break that here, and both turned out to be manageable. Markups that
move with size contribute a term whose sign is settled by a single condition --- no
group holding more than half of any market --- which is the same condition
Theorem~\ref{thm:entry} already needs. And a factory paying its parent's wage turns out
not to matter at all for the parent's country, because that country is paid a fixed
share of the group's whole bill, so the reshuffling of production inside the group
cancels. What is left, and what does not cancel, is a country whose income is
concentrated on one foreign partner.}

\noindent So the honest position is three sentences.

\begin{enumerate}
\item \textbf{The theorem.} Assumption~\ref{ass:H1} (no group above one half) and
Assumption~\ref{ass:H2} (the growth of the markets a country earns in exceeds its
exposure to the country whose wage moved) imply gross substitutes and hence at most one
equilibrium in the region where they hold.

\item \textbf{The conditions hold here, and are checked rather than assumed.} At the
calibrated economy the largest parent share anywhere is $0.34$ against a threshold of
$0.50$, and Assumption~\ref{ass:H2} holds for all $20$ ordered country pairs under both
ways of competing, with the tightest margin $0.024$ under quantity competition and
$0.025$ under price competition. Away from the solution, gross substitutes and the
contraction are certified on the \emph{continuum} of a box of wage vectors around the
$\nu=0$ solution by interval arithmetic --- every wage vector within $\pm15\%$ (log)
of the solved relative wages at once, with outward rounding --- and on a grid of the
wider region beyond it (spread $2.0$); the two reaches are reported separately.

\item \textbf{The second condition cannot be dropped, and the sufficient condition is
not the necessary one.} Proposition~\ref{prop:wagecounter} gives an economy in which
gross substitutes fails at the equilibrium. But in $50$ economies drawn from that same
family --- gross substitutes failing in $34$ of them --- every equilibrium still has
index $+1$, the index is independent of the numeraire in every case, and $25$ dispersed
starting points find one equilibrium in each. What uniqueness needs is the index
condition; gross substitutes is merely the tractable way to get it. Proving the index
condition directly is the remaining open problem, and it is a strictly smaller one than
the one this section used to describe.
\end{enumerate}

\noindent The old evidence still stands alongside the theorem, and it is the check that
the theorem cannot make, because the theorem holds the operating set fixed: ten random
starting wage vectors, spread over a range of roughly $e^{1.4}$, all reach wages
agreeing to $1.5\times10^{-11}$ with a labour-market error of $6.5\times10^{-12}$, and
all ten produce \textbf{the same set of operating firms}. Wages could in principle agree
while a different market structure prevailed, which is exactly the ambiguity
Theorem~\ref{thm:entry} rules out market by market.

\subsubsection{That an exact whole-number entry solution always exists}

Factories come in whole numbers, so an exact solution in whole numbers need not exist.
This is the classical integer problem of entry models, not something peculiar to this
one, and the honest response is to measure it. When no exact solution is reached, the
program keeps the arrangement leaving the fewest firms wanting to move, and reports the
remainder; it never breaks a tie by imposing an order. As the economy is enlarged:

\begin{center}
\begin{tabular}{rrrrr}
\toprule
firms & possible factory--market slots & operating & unresolved & share\\
\midrule
80  & 320  & 242 & 1 & 0.0031\\
160 & 640  & 369 & 0 & 0.0000\\
320 & 1280 & 582 & 3 & 0.0023\\
480 & 1920 & 768 & 2 & 0.0010\\
\bottomrule
\end{tabular}
\end{center}

\noindent The remainder is a fixed handful of borderline factories, so its share falls
roughly like $1/n$ and vanishes in a large economy. At the calibrated model it is 2
slots out of 1050 under quantity competition and 1 out of 1050 under price competition.

\begin{keypoint}
The distinction to insist on: this is a \textbf{near-miss on existence, not several
possible answers}. At no point does the model require choosing between two outcomes.
That is what Section~\ref{part:proofs} was for.
\end{keypoint}

\newpage
\section{Ownership and the Measurement of Foreign Profits}
\label{sec:result}
\label{sec:main}

\subsection{The result in words, before any algebra}

You want to know how much of the profit earned in a foreign market belongs to your
country. The natural calculation --- and the only one possible without firm-by-firm
ownership data --- is:
\[
  \text{(total profit earned there)}\times\text{(the fraction of those firms you own)}.
\]

This is wrong. Not merely imprecise: wrong in a direction you can predict, and by an
amount you can write down exactly.

The reason is Equation~\eqref{eq:profitquad}. Big firms are \emph{disproportionately}
profitable. So what matters is not what fraction of the \emph{firms} you own, but what
fraction of the \emph{big} ones. If the firms your country owns are larger than
average, the simple calculation \emph{understates} your share --- by exactly the
extent to which ownership and size go together.

\subsection{The result in symbols}

Let $\theta_g$ be the fraction of group $g$ owned by the country in question. In a
market with total spending $E$, that country's share of the profit is
\begin{equation}
  \frac{\Pi_H}{E}=\sum_g\theta_g\,S_g\,L(S_g).
  \label{eq:piH}
\end{equation}

\begin{proposition}[The correct adjustment uses concentration, not an average]
\label{prop:main}
Under quantity competition the margin is exactly linear, $L(S)=(1+\kappa S)/\sigma$,
and therefore
\begin{equation}
  \frac{\Pi_H}{E}
  =\frac{1}{\sigma}\Bigl[\;\underbrace{\bar\theta}_{\substack{\text{average}\\\text{ownership}}}
  \;+\;\kappa\bigl(\;\underbrace{\bar\theta\cdot \mathrm{HHI}}_{\substack{\text{concentration}\\\text{effect}}}
  \;+\;\underbrace{\mathrm{Cov}(\theta,S)}_{\substack{\text{ownership--size}\\\text{covariance}}}\bigr)\Bigr],
  \label{eq:decomp}
\end{equation}
where $\bar\theta=\sum_g S_g\theta_g$ is the sales-weighted ownership share,
$\mathrm{HHI}=\sum_g S_g^2$, and
$\mathrm{Cov}(\theta,S)=\sum_g\theta_gS_g^2-\bar\theta\cdot\mathrm{HHI}$.

A calculation using only the country-level ownership share computes $\bar\theta$ times
total profit, and therefore misses \emph{exactly}
\begin{equation}
  \boxed{\;\text{error}\;=\;\frac{E\,\kappa}{\sigma}\;\mathrm{Cov}(\theta,S).\;}
  \label{eq:error}
\end{equation}
\end{proposition}

\begin{proof}
Substituting $L(S)=(1+\kappa S)/\sigma$ into \eqref{eq:piH},
\[
  \frac{\Pi_H}{E}=\frac1\sigma\sum_g\theta_gS_g(1+\kappa S_g)
  =\frac1\sigma\Bigl[\sum_g\theta_gS_g+\kappa\sum_g\theta_gS_g^2\Bigr].
\]
The first sum is $\bar\theta$ by definition. For the second, the definition of
covariance gives $\sum_g\theta_gS_g^2=\bar\theta\cdot\mathrm{HHI}+\mathrm{Cov}(\theta,S)$,
which is \eqref{eq:decomp}. Total profit in the market is the case $\theta_g\equiv1$:
$\Pi/E=(1+\kappa\,\mathrm{HHI})/\sigma$, using $\sum_gS_g=1$. So the country-level
calculation is $\bar\theta\,\Pi/E=\bar\theta(1+\kappa\,\mathrm{HHI})/\sigma$.
Subtracting it from \eqref{eq:decomp}, the $\bar\theta$ and
$\kappa\bar\theta\,\mathrm{HHI}$ terms cancel, leaving exactly
$(\kappa/\sigma)\mathrm{Cov}(\theta,S)$.
\end{proof}

\says{Equation~\eqref{eq:decomp} splits a country's share of foreign profit into three
pieces: what it would get if all firms were the same size; a correction for the market
being concentrated; and a correction for \emph{which} firms it happens to own. The
country-level calculation captures the first two and misses the third entirely.
Equation~\eqref{eq:error} is that missing third piece. It is zero only if the firms a
country owns are no larger than average, or if $\kappa=0$.}

\why{This is the whole reason the project needed firm-level ownership data.
$\mathrm{Cov}(\theta,S)$ cannot be computed from a country-level ownership share --- it
requires knowing which \emph{particular} firms a country owns and how big each one is.
The prior literature listed in Section~\ref{sec:provenance} used country-level shares
throughout, so this term was invisible to it.}

\subsubsection{The object, measured in the data}
\label{sec:covmeasured}

The decomposition has now been computed in the matched customs--ownership data. A
market is one destination $\times$ product $\times$ year cell; a group is one ultimate
parent firm (the Orbis global-ultimate-owner identifier, with unmatched exporters as
their own singleton groups, the same convention as the six facts); shares are
within-LAC-sample shares, the same object Figure~6 measures. Aggregating markets with
value weights gives the three-term version of \eqref{eq:decomp}: the naive term
$\bar\theta\cdot\overline{\mathrm{HHI}}$ (computable from a country ownership share and
a published concentration index), a \emph{between-market} covariance (needs ownership
market by market), and a \emph{within-market} covariance (needs firm-level parent
identity --- this paper's data).

\begin{center}
\small
\begin{tabular}{lrrrrrr}
\toprule
owner & $\bar\theta$ & total & naive & between & within & total/naive\\
\midrule
United States & 0.135 & 0.0555 & 0.0503 & $+0.0032$ & $+0.0020$ & $\mathbf{1.10}$\\
Colombia      & 0.133 & 0.0614 & 0.0496 & $+0.0072$ & $+0.0046$ & 1.24\\
Chile         & 0.106 & 0.0434 & 0.0394 & $+0.0033$ & $+0.0006$ & 1.10\\
Argentina     & 0.096 & 0.0273 & 0.0357 & $-0.0072$ & $-0.0012$ & 0.76\\
United Kingdom& 0.088 & 0.0296 & 0.0328 & $-0.0035$ & $+0.0003$ & 0.90\\
Peru          & 0.067 & 0.0176 & 0.0251 & $-0.0058$ & $-0.0017$ & 0.70\\
\bottomrule
\end{tabular}
\end{center}

\noindent Three readings. \textbf{(i)} For the United States --- the tariff-setter of
the motivating question --- the country-level calculation \emph{understates} the
ownership-weighted concentration by $10\%$, and both covariance terms are positive:
US-owned firms are disproportionately the large players in the concentrated markets.
\textbf{(ii)} The correction is owner-specific and \emph{sign-varying}: Colombia's
owners are concentrated in exactly the markets they dominate ($+24\%$), while
Argentina's and Peru's sit in fragmented markets ($-24\%$, $-30\%$) --- a pattern no
country-level share can express, which is the contribution in one line.
\textbf{(iii)} The average market has measured $\mathrm{HHI}=0.37$ (within-sample
shares, mean $35$ groups per market), so the granular correction the proposition
multiplies by $\kappa/\sigma$ is far from the atomistic limit. The $\bar\theta$ column
independently reproduces the $\theta_{\mathrm{USA}}=0.134$ of the destination
tabulation, computed from a different cut of the same file.

\from{The data are the matched Orbis/D\&B--customs records behind the six facts (ten
Latin American origin countries, 2006--2022); the ultimate-parent identity is the
Orbis global-ultimate-owner identifier, with the conduit caveat of
Section~\ref{sec:fact1}'s data note applying here too. The three-term decomposition is
the market-aggregated form of Proposition~\ref{prop:main}'s identity; the
within-sample-share caveat, and the operator that maps such shares into
destination-absorption shares, are the measurement machinery introduced with Fact~4.
The computation is script 14 of the empirical repository.}

\begin{keypoint}
\eqref{eq:decomp} and \eqref{eq:error} are \textbf{identities}. Nothing is
approximated and no remainder is discarded. In the program they are checked against
the market solver's own profit figures and agree to $3\times10^{-16}$ and
$1\times10^{-12}$.

Under price competition the margin is not exactly linear, so \eqref{eq:decomp} becomes
the leading term of an expansion. The program then reports the exact value, the
approximation, and the gap, so the size of the approximation is measured rather than
assumed.
\end{keypoint}

\subsection{Why constant markups make the question disappear}

\begin{proposition}[Ownership irrelevance]\label{prop:irrelevance}
If all firms charge the same markup then $\kappa=0$, the second and third terms of
\eqref{eq:decomp} vanish, and a country's share of profit is exactly its average
ownership share times total profit: the distribution of ownership across firms is
irrelevant to \emph{income}, given the country-level share.
\end{proposition}

\begin{proof}
Immediate from \eqref{eq:decomp} with $\kappa=0$.
\end{proof}

\noindent(The statement is about the \emph{level} of profit income, which is what
\eqref{eq:decomp} proves. Whether the distribution also drops out of the \emph{optimal
tariff} at $\kappa=0$ involves how individual shares respond to the tariff,
$\sum_g\theta_g\,dS_g/dt$, and is not claimed here --- the same income-versus-policy
discipline as Section~\ref{sec:tariff}.)

\why{This proposition answers the reasonable objection ``why not use a simpler market
structure?''. In the simpler structure there is nothing to measure. It also explains
why the model cannot be built directly on the standard frameworks for multinational
production, attractive as their machinery is: those have constant or absent markups, so
the question of Section~1 is invisible inside them. This is not a criticism of that
work --- it is a statement about what a model must contain for this particular question
to be well posed.}

\subsection{Ownership and the optimal tariff: a result that runs against the original hypothesis}
\label{sec:tariff}

\begin{keypoint}
\textbf{This has to be stated plainly rather than buried.} The project began from the
idea that a country owning the foreign firms it taxes should want a \emph{lower}
tariff. In general equilibrium the model says the opposite.

\begin{center}
\begin{tabular}{lrr}
\toprule
Share of the foreign firms owned at home & Welfare effect of a small tariff & Optimal tariff\\
\midrule
0\%   & $-0.209$ & $<-0.90$\\
25\%  & $-0.146$ & $<-0.90$\\
50\%  & $-0.012$ & $0.00$\\
75\%  & $+0.190$ & $+0.20$\\
100\% & $+0.459$ & $+0.55$\\
\bottomrule
\end{tabular}
\end{center}

\smallskip
\noindent\emph{Status of this table: a directional preview, not a result.} It was
computed in the fixed-roster model, before endogenous entry, and its corner values
(optimal subsidies beyond $90\%$ at low ownership) are exactly the constant-marginal-cost
corner discussed in the text: with no upward-sloping foreign supply the markup
distortion dominates and optimal policy is a subsidy. Nothing in this table should be
quoted as a policy number; what it establishes is only the \emph{monotonicity} in
ownership, confirmed by two independent methods.
\end{keypoint}

\paragraph{The exercise redone with entry.} The full general equilibrium was re-solved
with the entry margin live, at every (tariff, ownership) pair: country 1 taxes all its
imports at rate $t$ and owns a share $\lambda$ of every foreign-headquartered parent.
Two findings, one per value of the input--output share.

\begin{center}
\small
\begin{tabular}{lccccc}
\toprule
$\nu=0$: welfare change vs.\ $t=0$ (\%) & $t=.05$ & $t=.10$ & $t=.20$ & $t=.30$ & optimal $t$\\
\midrule
$\lambda=0$    & $-0.06$ & $-0.23$ & $-0.97$ & $-1.81$ & $0$ (or subsidy)\\
$\lambda=0.50$ & $+0.15$ & $+0.22$ & $-0.10$ & $-0.58$ & $\approx0.10$\\
$\lambda=0.75$ & $+0.26$ & $+0.47$ & $+0.59$ & $+0.36$ & $\approx0.20$\\
$\lambda=1$    & $+0.69$ & $+1.05$ & $+1.38$ & $+1.48$ & $\ge0.30$\\
\bottomrule
\end{tabular}
\end{center}

\noindent At $\nu=0$ \textbf{the reversal survives entry cleanly}: the
welfare-maximising tariff rises monotonically with home ownership. At the calibrated
$\nu=0.55$, a tariff also taxes the imported input bundle, and ownership
\emph{neutralises} rather than flips the motive: the welfare cost of a $10\%$ tariff
shrinks from $-1.03\%$ at $\lambda=0$ to $-0.02\%$ at $\lambda=1$ --- a fifty-fold
compression of the gradient with the optimum still at (approximately) zero. The honest
statement is therefore: \emph{home ownership raises the optimal tariff monotonically;
whether it turns positive depends on the input--output structure.} The entry margin
itself is live and visible: tariffs drive affiliates out of the market entirely (from
$86$ to $73$ serving country 1 as $t$ runs to $30\%$ at $\lambda=0$), and exit varies
with ownership not through firm behaviour --- ownership never enters entry decisions
--- but through incomes: at $\lambda=1$ the taxing country is rich enough that its
market retains every affiliate at every tariff tested.

\from{The question is Brecher and Bhagwati's (1981) --- what does owning the taxed
foreign supplier do to the optimal tariff --- asked inside a model where the supplier
is an \emph{affiliate employing foreign labour} rather than an exporter, which is what
reverses the classical answer: a tariff depresses the host country's wage, and the
owner of the factories hiring that labour collects the saving. The grid design repeats
the fixed-roster preview above with the entry margin live; welfare is real income
$X_1/P_1$ with the Cobb--Douglas price index across sectors. Grid optima, not
first-order conditions: the ownership-adjusted sufficient-statistic formula for
$dW/dt$ in the entry model is still to be derived.}

\says{A tariff pushes down wages in the country being taxed. But the firms you own have
their \emph{factories} there, employing that now-cheaper labour. Their costs fall,
their profits rise --- and those profits are yours. Owning a foreign \emph{exporter} is
not the same as owning a foreign \emph{factory}. A tariff turns out to be partly a way
of lowering foreign wages, and if you own the capital employing those workers you
collect the gain twice: once as a consumer, once as a shareholder. The classical effect
runs the other way and is still present. It is outweighed here.}

\paragraph{What follows for how this is written up.}
\begin{itemize}
\item \textbf{Proposition~\ref{prop:main} stands regardless.} It is an accounting
identity, untouched by the direction of any policy experiment.
\item \textbf{Do not claim that foreign ownership lowers the optimal tariff.} It is not
established here and may be false. Two opposing forces operate and which dominates is
quantitative.
\item \textbf{The direction was an empirical question, and it has now been measured.}
It turns on where the observed factories actually sell. If they ship mostly back to the
parent's own country, the classical effect dominates; if they serve third countries
while employing local workers, the wage effect dominates. The tabulation is done: only
$9.2\%$ of foreign-owned export value ships to the parent's own country, so on average
the wage channel --- the uncomfortable one above --- is the empirically relevant case.
The origin split is large (El Salvador $0.40$, Dominican Republic $0.30$, Costa Rica
$0.28$, against Argentina, Chile and Peru near $0.06$), which makes ``when does home
ownership raise rather than lower the optimal tariff'' answerable in the cross-section.
\item \textbf{The exercise has now been redone with entry} (the table above): the
monotonicity survives, the exit channel is real and sized, and the input--output block
decides whether ownership merely neutralises the tariff's cost or turns the optimum
positive. Grid optima, not first-order conditions; the $\lambda=1$ column is a stark
thought experiment.
\end{itemize}

\why{There is a better paper in this than in the original hypothesis. ``When does
owning the firms you tax make you want to tax them \emph{more}, and why?'' is a sharper
question than re-deriving a known result with better data --- and it is the question
this model is equipped to answer.}

\newpage
\section{Quantification: Calibration and the Six Facts}
\label{sec:quant}

\subsection{Calibration}
\label{sec:calibration}

\begin{center}
\small
\begin{tabular}{llll}
\toprule
Parameter & Value & Status & Source, or the moment it is matched to\\
\midrule
$\sigma$, substitution within a sector & 5.0 & taken & trade elasticity: Simonovska--Waugh (2014) 4.14; Eaton--Kortum (2002) 8.28\\
$\eta$, substitution across sectors & 1.0 & taken & Atkeson--Burstein (2008): $\eta\approx1$, Cobb--Douglas\\
How firms compete & quantities & taken & prices supported and reported too\\
$\nu$, share of costs spent on inputs & 0.55 & taken & I--O shares as in Caliendo--Parro (2015); an economy-wide blend\\
$\omega$, own-sector share of inputs & 0.45 & taken & I--O tables are diagonal-heavy\\
$\alpha_k$, complexity & 0.10--0.55 & taken & head-office intensity, rising across sectors\\
$\gamma$, penalty for producing abroad & 1.18 & taken & below Head--Mayer's (2019) median MP friction of 31\%\\
Spread of firm ability & 4.0 & taken & Gaubert--Itskhoki (2021), Table 3: 4.38 at $\sigma=5$\\
Distance exponent & $1/(\sigma-1)$ & taken & makes the trade--distance elasticity $-1$\\
\midrule
$f$, cost of entering a market & 0.0006 & \textbf{matched} & the \emph{number of exporters}\\
Head-office disadvantage of local firms & 1.3 & \textbf{matched} & Fact 2, the gradient\\
$\zeta$, pool scaling with country size & 1.5 & \textbf{matched} & Fact 1, the foreign/domestic split\\
\midrule
Productivity edge of international firms & 0 & no longer needed & was matched to Fact 1's level; see below\\
Relative productivity of countries & 2.2 / 1.0 & calibrated & PWT 10 GDP per worker, advanced vs.\ LAC $\approx$ 2.3--2.7\\
Reference wage & 1.0 & units & unit of account (Walras)\\
\bottomrule
\end{tabular}
\end{center}

\begin{keypoint}
\textbf{The ``taken'' values have now been checked against their sources.} Three
notes survive the check and belong in the paper. (i) $\sigma=5$ is calibrated to the
\emph{trade-elasticity} literature (Simonovska--Waugh 4.14; Caliendo--Parro 4.5;
Eaton--Kortum 8.28), not to Atkeson--Burstein's own within-sector value of 10 --- markups
here are correspondingly larger than in their benchmark, which should be said.
(ii) $\nu=0.55$ is an economy-wide blend; manufacturing-specific input shares in
Caliendo--Parro's own input--output data run 0.6--0.7, so 0.55 is conservative.
(iii) $\gamma=1.18$ is conservative against Head--Mayer's median multinational-production
friction of 31\% in ad-valorem terms. The references in Section~\ref{sec:refs} have
been verified against the published versions.

So the degrees-of-freedom count is: \textbf{three internally matched parameters} ---
the entry cost, the head-office disadvantage, and the pool scaling --- against six
facts. A fourth, a multinational productivity edge matched to Fact 1's level, became
unnecessary when head-office services became a capability rather than a purchased
input, and is set to zero. Facts 3, 4 and 6 are therefore out-of-sample.
\end{keypoint}

\subsubsection{Adding entry pins down a parameter that used to float free}

Facts 2, 3 and 4 are all measured on the firms that \emph{export}. Without entry every
factory exported by construction, so that group was rich whatever the parameters, and
$f$ was not pinned down by anything. With entry it becomes an outcome --- and the first
guess for $f$ shrank the exporter group to six firms, at which point those three facts
cannot be measured at all.

\begin{center}
\begin{tabular}{rrrrl}
\toprule
$f$ & exporters & groups & home countries & concentration, factory $\to$ group\\
\midrule
0.0060 & 6  & 6  & 2 & $0.555\ (\times1.00)$\\
0.0020 & 31 & 24 & 3 & $0.148\ (\times1.40)$\\
\textbf{0.0006} & \textbf{59} & \textbf{43} & \textbf{4} & $\mathbf{0.080\ (\times1.51)}$\\
0.0002 & 71 & 53 & 4 & $0.071\ (\times1.51)$\\
\bottomrule
\end{tabular}
\end{center}

\says{An order of magnitude below the first guess, and flat below that. A parameter
that used to be arbitrary is now tied to something directly countable: how many firms
actually export.}

\subsection{Fact by fact}

\subsubsection{Fact 1: international firms are a large share of exports, and foreign ones dominate}
\label{sec:fact1}

\begin{center}
\begin{tabular}{lcc}
\toprule
Share of export value & model & data\\
\midrule
all international firms & 0.689 & 0.46--0.74\\
\quad foreign-owned & 0.543 & dominant in every country\\
\quad locally-owned & 0.146 & 0.02--0.28\\
\bottomrule
\end{tabular}
\end{center}

The \emph{level} is one matched parameter. The \emph{split} is not, and it is the more
informative half. It follows from the pool of possible parent firms growing with
country size (Section~\ref{sec:pool}): poor countries rarely draw a firm good enough to
become international. With pools of equal size everywhere the model puts locally-owned
international firms at $0.245$ --- \emph{larger} than foreign --- flatly contradicting
the data.

\subsubsection{Fact 2: foreign firms specialise in complex products}
\label{sec:fact2}

\begin{center}
\begin{tabular}{lcc}
\toprule
How the share changes with head-office intensity $\alpha_k$ & model & data\\
\midrule
foreign-owned international firms & $+0.42$ & $+0.43$\\
locally-owned international firms & $+0.58$ & negative\\
\bottomrule
\end{tabular}
\end{center}

\paragraph{The foreign half works, and required a structural change to get right.}
Without the head-office disadvantage the model gives $-0.58$: \emph{the wrong sign
entirely}. The reason is a genuine property of the cost-share \emph{variant} rather
than a coding error: there --- not in the baseline cost \eqref{eq:cost}, which has no
parent wage in it --- a foreign factory pays its \emph{parent's} wage on the $\alpha_k$
share of costs; parents sit in high-wage countries; so foreign ownership is most
expensive exactly where $\alpha_k$ is highest, and that variant predicts foreign shares
\emph{fall} with complexity. Something must overturn that, and the natural candidate
--- that head-office capability cannot be bought at arm's length, so a firm with no
parent is worse at precisely what complex products need --- is what the observed
gradient identifies. In the capability baseline the same disadvantage does the work
directly, and the gradient reaches $+0.42$ against the data's $+0.43$.

\paragraph{The locally-owned half does not work, and is reported rather than claimed.}
It comes out negative in a four-country world and positive in a five-country one, so the
model does not pin it down. The reason is transparent: a locally-owned international
firm is a local parent, so it escapes the head-office disadvantage --- and in the
cost-share variant it additionally pays a \emph{low} home wage on the $\alpha_k$ share,
the wrong-sign force again, now working in its favour. What the model does pin down is
that locally-owned international firms are \emph{few} (Fact 1), not what they
specialise in. This is an open problem.

\subsubsection{Fact 3: parent firms come from a handful of countries}

\begin{center}
\begin{tabular}{lcc}
\toprule
 & model & data\\
\midrule
concentration of parent countries & 0.273 & 0.130\\
largest single parent country & 0.344 & 0.247\\
\bottomrule
\end{tabular}
\end{center}

\textbf{Produced by the model, not assumed}, and that is the point. Nothing tells the
model which countries should host parent firms: every country draws from the same
distribution of ability, and only the number of draws differs. Concentration appears
because more draws means a better best draw. The model overshoots in a five-country
world, as it must, and earlier experiments show it moving towards the data as the world
is enlarged.

\medskip
\noindent\textbf{A caveat belonging to the data, not the model.} In the underlying
figure the largest parent country is the United Kingdom, ahead of the United States,
and several well-known conduit jurisdictions together account for roughly a tenth of the
total; about half of foreign-owned export value has no recorded parent at all. Any
ownership figure entering a table needs conduits reassigned to the true controlling
owner first.

\subsubsection{Fact 4: grouping factories by their owner raises measured concentration}

\begin{center}
\begin{tabular}{lccc}
\toprule
Concentration measured by & factory & group within a country & group worldwide\\
\midrule
model & 0.085 & 0.085 & 0.165 \quad($\times1.94$)\\
data  & 0.192 & 0.209 & 0.215 \quad($\times1.12$)\\
\bottomrule
\end{tabular}
\end{center}

\begin{keypoint}
\textbf{This is the fact Section~\ref{part:proofs} exists to protect.} The \emph{ordering} is a
prediction, not a matched moment. It requires markups to depend on the \emph{group's}
share --- cannibalisation --- and cannibalisation is exactly what the standard fix for
multiplicity cannot accommodate. Had the model taken that fix, this fact would have
been lost. Had it selected an outcome by assumption, market structure would have been
an artifact of the assumption. Theorem~\ref{thm:entry} is what lets the model keep the
mechanism \emph{and} make a determinate prediction.
\end{keypoint}

\noindent The size pattern also appears: groups operating in two or more countries are
15\% of groups and 72\% of export value, against 9\% of parents and 56\% of value for
the largest category in the data. The direction is right; the \emph{level} is not
comparable, because the model counts countries while the data count affiliates
worldwide.

\subsubsection{Fact 5: where there are more international firms, local firms export more}
\label{sec:fact5}

\begin{center}
\begin{tabular}{lc}
\toprule
Effect of adding international firms on local firms' exports, model & $-100\%$\\
The same experiment without endogenous entry & $-42\%$\\
Data & positive and statistically strong\\
\bottomrule
\end{tabular}
\end{center}

\begin{keypoint}
\textbf{This is the one clear failure, and adding entry made it worse.} The hope was
that international firms buying local inputs would push input prices down far enough to
let \emph{more} local firms clear their own entry threshold, flipping the sign for the
first time. Tested directly, it does not happen: international firms take business from
local ones on the entry margin too, and that outweighs the cheaper-inputs effect.
\end{keypoint}

Two explanations remained, and they could be told apart: a genuine spillover from
international to local firms' productivity, or contamination of the estimate --- the
controls in the published regression do not include a destination-by-product-by-year
term, so a demand shock in a particular market would raise international entry and
local exports together with no causal link.

\medskip
\noindent\textbf{The decisive test has now been run, and the fact survives.} Adding
destination$\times$product$\times$year fixed effects to the published specification's
tightest column (origin$\times$destination$\times$product and
origin$\times$destination$\times$year absorbed, clustering by origin--destination):

\begin{center}
\begin{tabular}{lcc}
\toprule
Effect of MNE presence on \emph{non-MNE} exports & published controls & + dest$\times$product$\times$year\\
\midrule
$\ln(\#\text{ MNE firms})$, intensive & $0.166$ $(0.014)$ & $\mathbf{0.087}$ $(0.017)$\\
any MNE present, extensive & $0.117$ $(0.011)$ & $\mathbf{0.061}$ $(0.011)$\\
PPML on the level, full controls & & $0.229$ $(0.051)$\\
\bottomrule
\end{tabular}
\end{center}

\noindent Roughly half of the published association was indeed common demand --- the
coefficient falls from $0.166$ to $0.087$ --- but the surviving half is precisely
estimated with the same sign on every margin. So Fact~5 is a \emph{real} target at
about half its published size, and the model's negative sign is a genuine failure that
can no longer be attributed to the specification.

\paragraph{What the mechanism is not, established by experiment.} Two spillover
channels have now been tested against the fact and both fail, in an instructive
pattern. A \emph{productivity} spillover (international presence raising local firms'
marginal productivity) moves the response only from $-97\%$ to $-60\%$ at sizes far
beyond anything estimated, because business stealing operates on the entry margin ---
an exited plant exports nothing however productive. A \emph{fixed-cost} spillover
(international presence lowering local plants' market-access cost) attacks that margin
directly, and at full strength it repairs it completely: every local plant--market
pair survives the influx of international firms. \textbf{The value response is still
$-80\%$.} With market spending fixed, entry is zero-sum in value, and pure
share-stealing on the intensive margin delivers the loss on its own. So no cost-side
spillover of any kind or size can produce the fact in this market structure.

\paragraph{What the mechanism must be --- and what happens when it is built.} The
data say the cell is \emph{not} zero-sum for the origin: total exports in a cell rise
strongly with international presence ($+0.69$ under the saturated controls), so the
origin's suppliers jointly capture more of the destination market. The model's markets
contained only the modelled countries' factories --- no rest-of-world supply inside a
market for the origin's firms to displace. The mechanism therefore has to work through
the origin's \emph{penetration}: an outside supplier in every market --- the $\lambda$
margin of the measurement operator, identified by the origin's exports over
destination absorption --- so that business stealing falls on the outside supply
rather than on fellow locals.

That outside supplier is now built (a rest-of-world country: large labour force, a
dense world-class fringe in every sector, no parents --- every uniqueness theorem
applies unchanged, and the largest group share \emph{falls}). Tested against the same
experiment, dose by dose:

\begin{center}
\begin{tabular}{lccc}
\toprule
 & in-sample share $\lambda$ & effect on local exports & local export pairs\\
\midrule
no outside supply & $1.00$ & $-97\%$ & $19\to1$\\
outside supply alone & $\approx0.28$ & $-71\%$ & $16\to7$\\
+ moderate spillovers & $\approx0.28$ & $-43\%$ & $33\to36$\\
same, $\lambda=0.11$ & $0.108$ & $+0.7\%$ & $29\to38$\\
\textbf{same, $\lambda=0.07$} & $\mathbf{0.071}$ & $\mathbf{+8.9\%}$ & $\mathbf{21\to38}$\\
\bottomrule
\end{tabular}
\end{center}

\noindent The margin is the first channel that moves the fact at first order on its
own; it flips the \emph{extensive} margin to the data's sign; and at the data's
$\lambda$ it flips the \emph{value} response too (the table's doses; exploratory,
with the spillover accruing to all locally-headquartered plants).

\paragraph{The adopted six-fact configuration.} Tuning on the full five-country world
adds one refinement: the spillovers accrue to \emph{single-plant local firms only} ---
the non-multinational locals that Fact~5 is about, and the population in which Aitken,
Hanson and Harrison measured the effect. Without that restriction, spillovers big
enough to carry Fact~5 also supercharge domestic multinationals' home plants and
invert Fact~1's ownership split. The adopted point is then: outside supplier at
$\lambda\approx0.07$, productivity spillover $0.17$, fixed-cost spillover $0.50$,
fringe-only --- \texttt{julia mne\_model.jl lambda}, committed output alongside. Its
scorecard, against the core calibration:

\begin{center}
\small
\begin{tabular}{lccl}
\toprule
 & core calibration & $\lambda$ configuration & data\\
\midrule
1: MNE share (foreign / domestic) & 0.69 (0.54/0.15) & 0.49 (0.39/0.11) & 0.46--0.74, foreign $\gg$\\
2: foreign gradient in complexity & $+0.42$ & $+0.58$ & $+0.43$\\
3: parent HHI / top & 0.27 / 0.34 & 0.29 / 0.42 & 0.13 / 0.25\\
4: grouping raises concentration & $\times1.94$ & $\times1.67$ & $\times1.12$\\
5: MNE presence $\to$ local exports & $-97\%$ & $\mathbf{+46\%}$ & positive ($+0.087$)\\
6: distance interactions & signs 3/3, too large & signs 2/3$^{*}$ & ---\\
\bottomrule
\end{tabular}
\end{center}

\noindent {\footnotesize $^{*}$In the $\lambda$ configuration the
already-present-at-destination interaction turns negative ($-0.39$ against the data's
$+0.07$); the other two distance coefficients keep their signs and move \emph{closer}
to the data than in the core. Disclosed, not hidden.}

\smallskip
\noindent Read the two columns as a trade, stated plainly: the core calibration
matches five facts with three fitted parameters and fails Fact~5; the $\lambda$
configuration turns Fact~5 \emph{positive on both margins} (magnitude within a factor
of three of the saturated estimate) and improves Fact~4, at the price of three more
fitted parameters, a Fact~2 overshoot, and one flipped sign inside Fact~6. Both are
reported so the reader sees exactly what the sixth fact costs. The remaining
discipline is to identify $\lambda$ from measured absorption shares (Comtrade plus
production data) rather than fitting it alongside Fact~5.

\from{The fact itself has a classic ancestor: Aitken, Hanson and Harrison (1997) found
that multinational presence raises neighbouring firms' probability of exporting, and
read it as export infrastructure --- exactly the fixed-cost channel tested above. The
productivity version is Javorcik (2004). Both mechanisms are implemented as options in
the program (\texttt{spill}, \texttt{fspill}), both computed off the roster of
potential plants so that the entry theorem is untouched, and in the six-fact
configuration their sizes are \emph{fitted} to Fact~5 rather than imported. The
outside supplier (\texttt{row\_L}) is the $\lambda$ of the measurement operator made
structural: the origin's exports over the destination's absorption, identifiable from
trade and production data (Comtrade imports plus domestic production) --- the
configuration's $\lambda=0.07$ sits inside the range those data imply for these
origins. It enters as an ordinary country --- labour, income, budget balance --- so
the accounting identities and every uniqueness theorem carry over without
modification.}

\subsubsection{Fact 6: distance matters less for international firms}

\begin{center}
\begin{tabular}{lcc}
\toprule
 & model & data\\
\midrule
effect of distance on exports & $-0.515$ & $-0.164$\\
$\quad$ smaller effect for international firms & $+0.701$ & $+0.046$\\
$\quad$ smaller still if already present at the destination & $+0.568$ & $+0.067$\\
\bottomrule
\end{tabular}
\end{center}

Signs and ordering are right, and the fact falls directly out of shipping costs being
counted from the factory rather than from head office --- point (iv) of
Equation~\eqref{eq:cost}. The magnitudes are too large. Note also that the estimate in
the data is an order of magnitude smaller than any standard gravity estimate, which
suggests a problem on the empirical side: the regression has no firm-level controls, so
the interaction is contaminated by which firms select into which markets.

\subsection{Scorecard}

\begin{center}
\begin{tabular}{clc}
\toprule
\# & Fact & Verdict\\
\midrule
1 & international firms a large share; foreign dominant & matched (level), \textbf{produced} (split)\\
2 & foreign firms in complex products & matched; \emph{local half unresolved}\\
3 & parents from a handful of countries & \textbf{produced}\\
4 & grouping by owner raises concentration & \textbf{produced} --- a prediction\\
5 & international firms raise local exports & fails in core ($-97\%$); \textbf{positive in the $\lambda$ configuration} ($+46\%$)\\
6 & distance matters less for international firms & \textbf{produced}, magnitudes too large\\
\bottomrule
\end{tabular}
\end{center}

Three parameters are matched in the core calibration: $f$ to the number of exporters,
the head-office disadvantage to Fact 2's gradient, and $\zeta$ to Fact 1's split.
There used to be a fourth --- a multinational productivity edge fitted to Fact~1's
level --- and it is no longer needed: once head-office services are a capability
rather than a factor bought at the parent's wage (Section~\ref{sec:hqchoice}),
multinationals are competitive enough without it and it is set to zero. Facts 3, 4
and 6 are therefore not fitted --- the model was not asked to reproduce them and does
so anyway. The six-fact $\lambda$ configuration adds three more (the outside
supplier's size and the two spillovers, matched to Fact~5 and the absorption share),
which is why it is reported as a named configuration beside the core rather than
silently replacing it: six facts against six fitted parameters is a different claim
from five facts against three, and the reader should see both.

\newpage
\section{Discussion: Limits and Next Steps}

\subsection{What this model does not do}

\begin{itemize}
\item \textbf{Fact 5 needs the outside supplier.} The fact survives the saturated
fixed-effects test at half its published size, the \emph{core} calibration has the
wrong sign, and cost-side spillovers alone cannot fix it. The six-fact $\lambda$
configuration (Section~\ref{sec:fact5}) turns it positive on both margins, with the
magnitude within a factor of three --- at the price of three additional fitted
parameters, a Fact~2 overshoot, one flipped sign inside Fact~6, and an outside
supplier whose size should ultimately be disciplined by measured absorption shares
rather than by Fact 5 itself. Both configurations are reported so the reader sees
exactly what the fact costs.
\item \textbf{The locally-owned half of Fact 2 is not pinned down}
(Section~\ref{sec:fact2}). The model determines that such firms are few, not what they
specialise in.
\item \textbf{Uniqueness of wages rests on a condition that is checked, not on nothing.} The strongest statement is Theorem~\ref{thm:contract}: the map the solver iterates is a contraction in Hilbert's metric wherever its Jacobian is entrywise positive, so the model has one solution and the algorithm finds it. That positivity is gross substitutes plus a damping bound the model itself supplies. In the $\nu=0$ version it is now \emph{computer-certified on the continuum} of a box of wage vectors around the solution (outward-rounded interval arithmetic --- no between-points gap), and beyond that box it holds at every point tested, out to spreads of $2.0$, under both ways of competing. What is not available is an \emph{unconditional} theorem, and the reason is stated in Section~\ref{sec:notproved}: the last channel that could break the condition is the ordinary income effect, which is this paper's own subject. Proposition~\ref{prop:wagecounter} shows the condition can fail in the variant where head-office services are bought at the parent's wage, which is one reason the baseline is not that variant.
Uniqueness within a market and of the entry outcome \emph{are} proved.
\item \textbf{Ownership $\theta$ is taken as given.} Nobody chooses a portfolio. Making
ownership a choice, alongside large firms and endogenous policy, is a substantially
larger project.
\item \textbf{Factory locations are candidates, not draws from a continuum.} Each group
has one candidate factory per country and entry selects among them. The richer version,
in which a firm's ability is correlated across its locations, is available as a switch
but is not estimated, and the literature disagrees about its value.
\item \textbf{The tariff exercise now exists with entry} (Section~\ref{sec:tariff}),
as grid optima rather than first-order conditions; the ownership-adjusted
sufficient-statistic formula for $dW/dt$ remains to be derived in the entry model.
\item \textbf{Parameter values marked ``taken'' have not been re-verified} against their
sources for this document. The references have.
\end{itemize}

\subsection{What to do next, in order --- and what is already done}

\begin{enumerate}
\item \textbf{Where multinational exports go is now measured}, and it settles the
direction of Section~\ref{sec:tariff}'s question: only $9.2\%$ of foreign-owned export
value ships to the parent's own country, foreign affiliates are \emph{less} US-oriented
than local exporters, and the share of US imports from these origins produced by
US-owned affiliates is $0.134$. The classical rent-repatriation channel is weak on
average --- but it varies enormously by origin (from $0.40$ in El Salvador to $0.02$
in Paraguay), which turns the policy question into a cross-sectional one.
\item \textbf{Conduit jurisdictions have been reassigned} to the first non-conduit
controlling owner, and the headline barely moves: the US ownership share goes
$0.134\to0.136$, with an unresolvable residual disclosed as the range $[0.136,0.170]$.
\item \textbf{Fact 5 is settled in two steps.} Re-estimated with
destination-by-product-by-year controls, it survives at half size ($0.166\to0.087$
intensive, $0.117\to0.061$ extensive, PPML agreeing); and the six-fact $\lambda$
configuration now turns it positive on both margins ($+46\%$ at the experiment's dose;
Section~\ref{sec:fact5}). What remains is discipline, not discovery: calibrate the
outside supplier to measured absorption shares (Comtrade plus production data) so that
$\lambda$ is identified independently rather than fitted alongside Fact 5.
\item \textbf{Measure the intra-group share of those exports} (needs the affiliate
roster, not the customs file alone).
\item \textbf{The tariff exercise with entry is done} (Section~\ref{sec:tariff}):
monotonicity in ownership survives; the sign of the optimum depends on the
input--output block. Next on the policy side: the ownership-adjusted
sufficient-statistic formula for $dW/dt$ in the entry model.
\end{enumerate}

\newpage
\section*{References}
\label{sec:refs}

Years and outlets checked against the published sources.

\begin{itemize}[leftmargin=1.4em]
\item Aitken, B., G. Hanson and A. Harrison (1997). ``Spillovers, Foreign Investment,
and Export Behavior.'' \emph{Journal of International Economics} 43(1--2), 103--132.
\item Antr\`as, P. (2003). ``Firms, Contracts, and Trade Structure.''
\emph{Quarterly Journal of Economics} 118(4), 1375--1418.
\item Antr\`as, P. and E. Helpman (2004). ``Global Sourcing.''
\emph{Journal of Political Economy} 112(3), 552--580.
\item Alviarez, V., K. Head and T. Mayer (2025). ``Global Giants and Local Stars: How
Changes in Brand Ownership Affect Competition.'' \emph{American Economic Journal:
Microeconomics} 17(1), 389--432.
\item Atkeson, A. and A. Burstein (2008). ``Pricing-to-Market, Trade Costs, and
International Relative Prices.'' \emph{American Economic Review} 98(5), 1998--2031.
\item Birkhoff, G. (1957). ``Extensions of Jentzsch's Theorem.''
\emph{Transactions of the American Mathematical Society} 85, 219--227.
\item Blanchard, E. (2007). ``Foreign Direct Investment, Endogenous Tariffs, and
Preferential Trade Agreements.'' \emph{B.E. Journal of Economic Analysis \& Policy}
7(1).
\item Blanchard, E. (2010). ``Reevaluating the Role of Trade Agreements: Does
Investment Globalization Make the WTO Obsolete?'' \emph{Journal of International
Economics} 82(1), 63--72.
\item Blanchard, E. and X. Matschke (2015). ``U.S. Multinationals and Preferential
Market Access.'' \emph{Review of Economics and Statistics} 97(4), 839--854.
\item Blanchard, E., C. Bown and R. Johnson (2026). ``Global Value Chains and Trade
Policy.'' \emph{Review of Economic Studies} 93(1), 181--214.
\item Brecher, R. and J. Bhagwati (1981). ``Foreign Ownership and the Theory of Trade
and Welfare.'' \emph{Journal of Political Economy} 89(3), 497--511.
\item Caliendo, L. and F. Parro (2015). ``Estimates of the Trade and Welfare Effects of
NAFTA.'' \emph{Review of Economic Studies} 82(1), 1--44.
\item Eaton, J. and S. Kortum (2002). ``Technology, Geography, and Trade.''
\emph{Econometrica} 70(5), 1741--1779.
\item Gaubert, C. and O. Itskhoki (2021). ``Granular Comparative Advantage.''
\emph{Journal of Political Economy} 129(3), 871--939.
\item Gaubert, C., O. Itskhoki and M. Vogler (2021). ``Government Policies in a
Granular Global Economy.'' \emph{Journal of Monetary Economics} 121, 95--112.
\item Head, K. and T. Mayer (2019). ``Brands in Motion: How Frictions Shape
Multinational Production.'' \emph{American Economic Review} 109(9), 3073--3124.
\item Itskhoki, O. and D. Mukhin (2025). ``The Optimal Macro Tariff.'' NBER Working
Paper 33839.
\item Javorcik, B. (2004). ``Does Foreign Direct Investment Increase the Productivity
of Domestic Firms? In Search of Spillovers Through Backward Linkages.''
\emph{American Economic Review} 94(3), 605--627.
\item Milgrom, P. and C. Shannon (1994). ``Monotone Comparative Statics.''
\emph{Econometrica} 62(1), 157--180.
\item Moore, R. E. (1966). \emph{Interval Analysis}. Prentice-Hall.
\item Novshek, W. (1985). ``On the Existence of Cournot Equilibrium.''
\emph{Review of Economic Studies} 52(1), 85--98.
\item Ramondo, N. and A. Rodr\'iguez-Clare (2013). ``Trade, Multinational Production,
and the Gains from Openness.'' \emph{Journal of Political Economy} 121(2), 273--322.
\item Simonovska, I. and M. Waugh (2014). ``The Elasticity of Trade: Estimates and
Evidence.'' \emph{Journal of International Economics} 92(1), 34--50.
\item Tintelnot, F. (2017). ``Global Production with Export Platforms.''
\emph{Quarterly Journal of Economics} 132(1), 157--209.
\item Yang, C. (2023). ``Location Choices of Multi-plant Oligopolists: Theory and
Evidence from the Cement Industry.'' Working paper (R\&R, \emph{JPE: Macroeconomics}).
\end{itemize}

\newpage
\appendix
\part*{Appendix}

\section{Notation: every symbol used in this document}
\label{sec:notation}

\begin{center}
\small
\begin{longtable}{ll p{8.4cm}}
\toprule
\textbf{Symbol} & \textbf{Name} & \textbf{Meaning}\\
\midrule
\endfirsthead
\toprule
\textbf{Symbol} & \textbf{Name} & \textbf{Meaning}\\
\midrule
\endhead
\multicolumn{3}{l}{\emph{Indices}}\\
$n$ & destination & the country where a good is sold\\
$\ell$ & host & the country where a good is produced\\
$h$ & home & the country where the owning group has its head office\\
$k$ & sector & one product category\\
$i$ or $a$ & factory & one production plant, in one country, in one sector\\
$g$ & group & the firm that ultimately owns a set of factories\\
\midrule
\multicolumn{3}{l}{\emph{Demand}}\\
$\sigma_k$ & inner elasticity & how readily buyers switch between goods \emph{within} sector $k$\\
$\eta$ & outer elasticity & how readily buyers switch \emph{away from} a sector altogether; $\sigma_k>\eta\ge1$\\
$\beta_k$ & taste weight & share of final spending going to sector $k$\\
$E$ & market size & total spending in one (destination, sector) market\\
$P_{nk}$ & price index & cost of one unit of sector $k$'s bundle in country $n$\\
\midrule
\multicolumn{3}{l}{\emph{Competition}}\\
$S_g$ & group share & group $g$'s share of sales in a market, adding up all its factories\\
$\mu(S)$ & markup & price divided by marginal cost; rises with $S_g$\\
$L(S)$ & margin & profit per unit of sales, $L=1-1/\mu$\\
$\kappa$ & profit curvature & how sharply profit rises with size; the central result is proportional to it\\
$\Pi_g$ & profit & group $g$'s profit in a market, before entry costs\\
$\mathrm{HHI}$ & concentration & $\sum_g S_g^2$, the standard concentration index\\
\midrule
\multicolumn{3}{l}{\emph{The two objects that make the model tractable}}\\
$K_g$ & \textbf{strength} & $\sum_{i\in g}c_i^{1-\sigma}$: group $g$'s entire position in one number\\
$A$ & market index & $\sum_i p_i^{1-\sigma}$: all a group needs to know about its rivals\\
$\psi(S)$ & share map & $S\mu(S)^{\sigma-1}$; converts strength into share, $S_g=\psi^{-1}(K_g/A)$\\
$\Omega(K;A)$ & entry payoff & a group's profit written as a function of its own strength\\
\midrule
\multicolumn{3}{l}{\emph{Costs and technology}}\\
$c_{a,n}$ & delivered cost & cost of making one unit at factory $a$ and getting it to $n$\\
$\varphi_a$ & productivity & how efficient factory $a$ is\\
$w_n$ & wage & wage in country $n$; one country's wage is the unit of account\\
$\alpha_k$ & complexity & head-office intensity of sector $k$; enters productivity, not cost\\
$\nu_k$ & input share & fraction of costs spent on goods bought from other firms\\
$\omega_{k'k}$ & input mix & share of sector $k$'s input bill spent on sector $k'$; sums to one over $k'$\\
$\gamma_{h\ell}$ & distance-from-HQ penalty & efficiency loss from producing away from head office\\
$d_{\ell n}$ & shipping cost & cost of moving goods from $\ell$ to $n$\\
$t_{\ell n k}$ & tariff & tariff on sector $k$ goods moving from $\ell$ to $n$\\
$F_{a,n}$ & entry cost & fixed cost for factory $a$ to sell into market $n$\\
$f_k$ & entry-cost scale & the size of that fixed cost in sector $k$\\
\midrule
\multicolumn{3}{l}{\emph{Countries and ownership}}\\
$L_n$ & workforce & number of workers in country $n$, all of whom must be employed\\
$X_n$ & income & total income of country $n$\\
$\theta_{gn}$ & ownership & fraction of group $g$ owned by country $n$; sums to one over $n$\\
$\bar\theta$ & average ownership & sales-weighted average of $\theta_g$ in a market\\
$\zeta$ & pool scaling & how fast the number of possible groups grows with country size\\
\bottomrule
\end{longtable}
\end{center}

\addcontentsline{toc}{part}{Appendix}

\section{Summary of the model's equations}

For reference, the model in one page. Every symbol is defined in
Section~\ref{sec:notation}.

\begin{center}
\begin{tabular}{clp{7.2cm}}
\toprule
Eq. & Statement & What it determines\\
\midrule
\eqref{eq:muC}--\eqref{eq:muB} & $\mu(S_g)$ & the markup a group charges, rising with its own share\\
\eqref{eq:profit} & $\Pi_g=E\,S_g\,L(S_g)$ & the group's profit before entry costs\\
\eqref{eq:clearing} & $\sum_g S_g=1$ & pins down the market index $A$\\
\eqref{eq:psiinv} & $S_g=\psi^{-1}(K_g/A)$ & each group's share, from its strength and market toughness\\
\eqref{eq:cost} & $c_{a,n}$ & what one unit costs to make and deliver\\
\eqref{eq:fixed} & $F_{a,n}$ & the fixed cost of selling into a market; determines who enters\\
\eqref{eq:income} & $X_n$ & each country's income, including profits it owns\\
\eqref{eq:decomp} & $\Pi_H/E$ & the share of foreign profit accruing to one country\\
\eqref{eq:error} & $(E\kappa/\sigma)\mathrm{Cov}(\theta,S)$ & what a country-level ownership share misses\\
\bottomrule
\end{tabular}
\end{center}

\section{Reproducing everything}

One program, nothing to install beyond the language itself: \texttt{mne\_model.jl}
holds the whole model, and \texttt{julia mne\_model.jl <mode>} reproduces the
correspondingly named part of this document.

\begin{center}
\begin{tabular}{ll}
\toprule
\texttt{core} & one market: Theorem~\ref{thm:market} and Proposition~\ref{prop:main}\\
\texttt{entry} & Theorem~\ref{thm:entry}: the closed forms and the brute-force check\\
\texttt{ge} & the accounting checks, existence, and the uniqueness protocol\\
\texttt{facts} & the six facts of Section~\ref{sec:quant}\\
\texttt{bertrand} & all of the above under price competition\\
\texttt{quick} & everything on a smaller world, in a few minutes\\
\bottomrule
\end{tabular}
\end{center}

Three companion programs, each with its committed output alongside:
\texttt{wage\_uniqueness.jl} (Theorem~\ref{thm:wage} and the counterexample),
\texttt{simple\_model.jl} (Theorem~\ref{thm:simple} and the baseline comparison),
\texttt{interval\_certificate.jl} (the continuum certificate: outward-rounded
interval arithmetic over whole boxes of wage vectors), and
\texttt{experiments.jl} (the one-off parameter scans quoted in the text ---
the spillover grids, the $\nu$ comparison, and the tariff-with-entry exercise).

\end{document}
